
\section{Introduction}
In this chapter, we present  computational experiments to empirically
validate the results obtained in previous three chapters and demonstrate the
performance of the algorithms used to solve LRSP. In this chapter, first we
describe the details of the instances used in our experiments. Second, we
present computational results related to Chapter 1. In this context, we compare
the LP relaxation bounds for the LRSP formulations and validate the effect of
valid inequalities presented in the first chapter. Third, we get together a set of
experiments to demonstrate the evolution of the branch and price algorithms for
LRSP described in the second chapter. This section evaluates the components of
the branch and price algorithms as well as the performance of the whole algorithms.
Finally, the last section displays the effect of extending the algorithm to a branch
price and cut algorithm using the cutting planes described in the third chapter.
We investigate each implemented cutting plane.

As noted in Chapter 3, the branch, price and cut algorithms were coded in MINTO with
CPLEX 9.1 as an LP solver. Unless otherwise stated for the experiments, we implemented
the experiments in MINTO.  We ran our experiments on a
64-node Linux-based Beowulf cluster using condor.
%For all of the experiments, we used a Linux-based workstation
%with a 1.8 GHz processor and 2GB RAM. We ran our experiments using condor.


\section{Description of the Instances}
\label{instances}

The only LRSP instances available in the literature are those of \citet{Lin}.
We could get only four small instances from \citet{Lin}. We describe those instances
and present our results for those  in Section~\ref{compareLin}.
Since they are very limited and small in size, for testing, we generated several sets
of new instances.

\paragraph{Location and Demand Data}
We generated four sets of base instances from the MDVRP instances of \citet{Cordeau}:
instances with 20, 25, 30 and 40 customers. Each instance contains a subset of the
customer locations and demand data from one of the \citet{Cordeau} instances.
In particular, we selected instances $p01$ (50 customers), $p03$ (75 customers),
$p07$ (100 customers) and $p11$ (249 customers) from which to generate our base sets
of instances (these reference instances differ from each other in the locations
of the customers and the candidate facilities).
%These three instances
%were selected because they differ from each other in the locations of the customers and the
%candidate facilities.
%\paragraph{Location Data}
%For our main experiments, we generated two sets of eight base instances;
%in one set, each instance has 25 customers and in the other, each instance
%has 40 customers.

Table \ref{Locationdata} presents the set of customers used to construct LRSP
instances in each group of instances (20, 25, 30 and 40 customer instances).
Table \ref{Locationdata} specifies the marker letter used to label our instances,
the corresponding \citet{Cordeau} instance (denoted as {\em Ref}) and the indices
of customers included from the reference instance.
% for each instance in each instance group . The marker letter was used to label the
%LRSP instances.
The details of labeling is explained later in this Section.
In addition to the customer indices, Table \ref{FacLocdata} lists the coordinates of the
five facility locations in each instance group differentiated with each reference
%\citet{Cordeau}
instance.
% and the coordinates of the five
%candidate facility locations
%(denoted as {\em Fac. Loc. (X,Y)}).
%The instances generated from the same \citet{Cordeau} instance have the same set
%of facility locations.
\begin {table}[h]
\begin{center}
\caption {Location \& Demand Data for LRSP instances generated from \citet{Cordeau}} \label{Locationdata}
\begin {tabular}[h]{|c|c|c|c|c|c|} \hline
Marker  & \mbox{} & \multicolumn{4}{|c|}{Indices of the Coordinates from Reference} \\\cline{3-6}
%Marker & \mbox{} & \multicolumn{4}{|c|}{Instance Group} \\\cline{3-6}
Letter  & Ref     & 20-cust. &  25-cust. &  30-cust.&  40-cust. \\\hline
f & p01 & 1 - 20  & 1 - 25   & 1-30 & 1-40       \\
l & p01 & 31 - 50 & 26 - 50  & -    & 11-50      \\\hline
f & p03 & 1 - 20  & 1 - 25   & 1-30 & 1-40        \\
m & p03 & 28 - 47 & 26 - 50  & -    & 16-35, 41-60 \\
l & p03 & 56 - 75 & 51 - 75  & 46-75    & 36-75       \\\hline
f & p07 & 1 - 20  & 1 - 25   & 1-30 & 1-40        \\
s & p07 & 21 - 40 & 26 - 50  & -    & 31-70       \\
t & p07 & 31 - 50 & 51 - 75  & -    & 1-20, 51-70 \\\hline
f & p11 & -   & -        & 1-30 & -           \\
l & p11 & -   & -        & 101-130  & -           \\\hline
\end{tabular}
\end{center}
\end{table}

\begin {table}[h]
\begin{center}
\caption {Facility Locations for LRSP instances} \label{FacLocdata}
\begin {tabular}[h]{|l|l|l|} \hline
Instance Group & Ref & Facility Coordinates (X,Y) \\\hline
20- and 25-customer   & p01 &   (20,20), (30,40), (15,40), (45,40) (55,50) \\
\mbox{} &   p03   & (40,40), (50,22),(25,45), (20,20),  (55,55)  \\
\mbox{} &   p07   & (15,35), (55,35),(35,20), (35,50), (25,45)   \\\hline
30-customer &   p01   & (25,20), (40,40),(50,30), (55,55), (15,45)   \\
\mbox{} &   p03   & (20,20), (40,40),(50,22), (55,55), (25,45)    \\
\mbox{} &   p07   & (15,35), (55,35),(35,20), (35,50), (25,45)     \\
\mbox{} &   p11-f & (70,0), (40,80),(40,-80), (-60,20), (-60,-20) \\
\mbox{} &   p11-l & (70,0), (40,80),(20,-30), (-60,20), (-60,-20) \\\hline
40-customer &   p01   & (20,20), (30,40), (15,40), (50, 30), (55,55)  \\
\mbox{} &   p03   & (40,40), (50,22), (25,45), (20,20), (55,55)  \\
\mbox{} &   p07   & (15,35), (55,35),(35,20), (35,50), (25,45)  \\\hline
\end{tabular}
\end{center}
\end{table}

\begin{comment}
\begin {table}[h]
\begin{center}
\caption {Details for LRSP instances generated from \citet{Cordeau} MDVRP instances} \label{Locationdata}
\begin {tabular}[t]{|c|c|c|c||c|c|c|c|} \hline
\# & Ref & Customers & Fac. Loc. (X,Y) & \# & Ref & Customers  & Fac. Loc. (X,Y) \\ \hline
1 & p01 & 1 - 25 & (20,20), (30,40), (15,40), & 9 & p01 & 1 - 40 & (20,20), (30,40), (15,40),\\
2 & p01 & 26 - 50 & (45, 40), (55,50) & 10 & p01 & 11 - 50 & (50, 30), (55,55) \\ \hline
3 & p03 & 1 - 25 & (40,40), (50,22),  & 11 & p03 & 1 - 40 & (40,40), (50,22) \\
4 & p03 & 26 - 50 & (25,45), (20,20), & 12 & p03 & 16 - 35, 41 - 60 & (25,45), (20,20) \\
5 & p03 & 51 - 75 & (55,55) & 13 & p03 & 36 - 75 & (55,55) \\ \hline
6 & p07 & 1 - 25 & (15,35), (55,35), & 14 & p07 & 1 - 40 & (15,35), (55,35), \\
7 & p07 & 26 - 50 & (35,20), (35,50), & 15 & p07 & 31 - 70 & (35,20), (35,50)\\
8 & p07 & 51 - 75 &  (25,45) & 16 & p07 & 1 - 20, 50 - 70 &  (25,45)\\ \hline
\end{tabular}
\end{center}
\end{table}
\begin {table}[h]
\begin{center}
\caption {Details for LRSP instances generated from \citet{Cordeau} MDVRP instances} \label{Locationdata}
\begin {tabular}[t]{|c|c|c||c|c|c|} \hline
\# & Ref & Customers & \# & Ref & Customers \\ \hline
1 & p01 & 1 - 25 & 9 & p01 & 1 - 40 \\
2 & p01 & 26 - 50 & 10 & p01 & 11 - 50 \\ \hline
3 & p03 & 1 - 25 & 11 & p03 & 1 - 40 \\
4 & p03 & 26 - 50 & 12 & p03 & 16 - 35, 41 - 60 \\
5 & p03 & 51 - 75 & 13 & p03 & 36 - 75 \\ \hline
6 & p07 & 1 - 25 & 14 & p07 & 1 - 40 \\
7 & p07 & 26 - 50 & 15 & p07 & 31 - 70 \\
8 & p07 & 51 - 75 & 16 & p07 & 1 - 20, 50 - 70 \\ \hline
\end{tabular}
\end{center}
\end{table}
\end{comment}

\paragraph{Distance Matrix}
Rounded Euclidean distances between each pair of nodes were used in
the algorithm, and we ensured that the distance matrices satisfied the triangle
inequality.
\paragraph{Capacity and Cost Parameters}
%We used the demands from \citet{Cordeau} corresponding to the customers
%included in each instance.
Based on the demand information, we assumed that the capacity of all facilities
was 250, which requires at least two facilities to be open for each instance.

%We expected the performance of the branch-and-price algorithm to be most
%affected by the time limit and the vehicle capacity values. Therefore,
For each instance, we considered two time limit values and two vehicle capacity
values. Time limits of 7 and 8 hours (which can be converted to distance limits of
140 and 160 miles under the assumption that vehicles travel 20 miles per hour) were used.
In all instances, 140 is more than twice the distance between each customer and its
farthest candidate facility. To be able to satisfy this condition, for instances
generated from $p11$, we used distance limits of 250 and 270 instead of 140 and 160.
We used ``t1'' or ``t2'' to mark the instances based on distance limit; ``t1'' was
used for instances with distance limit 140 (or 240) and ``t2'' was used for instances with
distance limit 160 (or 260).

We chose two vehicle capacity values. To determine these values, we used
the fact that at least two facilities are required to be open (due to the capacity)
and used the following equation similar to \citet{Laporte86}:
\begin{equation}
\label{veccap-calc}
B = (1-\alpha) \mbox{max}_{i \in I}\{D_i\} + \alpha \frac{\sum_{i \in I}D_i}{2},
\end{equation}
with $\alpha$ equal to 0.1 and 0.2. Small $\alpha$ values were considered since
a schedule employing multiple trips for a single vehicle is only realistic
if the vehicle capacity is small compared to the demands of the customers.
We used the 25 customer instances to determine vehicle capacities and used
the same values in the 20, 30 and 40 customer instances.
For each $\alpha$ value, we took the average of the calculated $B$ values
for the instances generated from the same \citet{Cordeau} instance, divided it by 10,
rounded it up, and multiplied by 10.
We obtained vehicle capacity values 50 (for $\alpha=0.1$) and  70 (for $\alpha=0.2$)
for instances generated from $p01$ and $p07$, 60 (for $\alpha=0.1$) and 80
(for $\alpha=0.2$) for the instances generated from $p03$, and 150 and 200
for the instances generated from $p11$. The smaller vehicle capacity
values can cover at least one customer plus the customer that has the largest
demand in the instance. We used ``v1'' or ``v2'' to mark the instances based on
vehicle capacity; ``v1'' was used to represent smaller capacity value
and ``v2'' was used for larger capacity value.

Based on the cost values reported by \citet{Burlett} and  Kenworth Truck Company
\citeyearpar{Kenworth}, fixed costs for the facilities were generated
using a uniform distribution on the interval [1500,2300],
the fixed cost and the operating cost of a mid-size vehicle were estimated to be
\$225 per day and \$1 per 1 mile traveled, respectively.
The same cost parameters were used in all instances.

Table \ref{probparam} summarizes the capacity and cost parameters. In the Table,
as in Chapter 1, $L^V$ and $L^T$ represent the vehicle capacity and distance limit,
respectively. $C^F_j$ denotes the fixed cost of the facility $j$ for all $j \in J$
while $C^V$ and $C^O$ correspond the vehicle fixed and operating costs, respectively.
\begin {table}[h]
\begin{center}
\caption {Summary of Parameters} \label{probparam}
\begin {tabular}[t]{lccccc}
%\multicolumn {6}{c}{{\bf Table 1.} Problem Parameters}\\ \hline
Generated Using & $L^V_1$   & $L^V_2$ & $L^T_1$ & $L^T_2$ & Fac. cap.\\\hline
p01,p07 &   50  &   70  &   140 &   160 &   250 \\
p03 &   60  &   80  &   140 &   160 &   250\\
p11 &   150 &   200 &   240 &   260 &   250\\\hline
\\
\end{tabular}
\begin {tabular}[t]{ccccccc}
$C^F_1$ & $C^F_2$ & $C^F_3$ &$C^F_4$ & $C^F_5$ & $C^V$ & $C^O$ \\\hline
\$1615 /day  &  \$1845 /day  &  \$1781 /day  &  \$1975 /day  &  \$1753 /day  &  \$225 /day  & \$ 1 /mile \\\hline
\end{tabular}
\end{center}
\end{table}

\paragraph{Labels of the LRSP Instances}
We label our instances considering reference \citet{Cordeau} instances,
set of customers included from the corresponding reference
(the marker letters listed in Table \ref{Locationdata} used to indicate this
information), the number of customers, vehicle capacity (one of the markers
v1 or v2) and distance limit values (one of t1 or t2).
%based on the reference \citet{Cordeau} instances,
%the marker letters showing the set of customer coordinates used (according
%to Table \ref{Locationdata}), the number of customers, vehicle capacity and
%distance limit values.
For example, we used label ``p01-f25-v1t1'' for the
LRSP instance (in 25-customer instance group) generated using customers 1 to
25 (based on letter ``f'') in $p01$. The ``v1'' and ``t1'' markers inform that
the smaller vehicle capacity and distance limit values (50 and 140, respectively,
based on Table \ref{probparam}) are used in the instance. Another example instance
label is  ``p03-m40-v1t2''. This label denotes that the instances is one of the
40-customer instances and is generated from reference $p03$ using customers 16 to 35,
41 to 60. Vehicle capacity and distance limits are 60 and 160, respectively.

\section{Computational Experiments For Chapter 1}
The goal of the computational experiments in this section
is two-fold: to compare the LP relaxation bounds of G-LRSP and SPP-LRSP
and to empirically demonstrate the benefit of the valid inequalities described
in Section \ref{StrongForm}.

\subsection{Comparing the LP Relaxations of SPP-LRSP and G-LRSP}

In Section \ref{part1:dantzigW}, it is shown that the Dantzig Wolfe
decomposition of graph-based model G-LRSP is equal to the set
partitioning based model, SPP-LRSP. In Section \ref{sec:comp-form}
we give the theoretical comparison between the LP relaxation
bounds of G-LRSP and SPP-LRSP. In this section, we validate this
result and empirically demonstrate the strength of the LP relaxation of
SPP-LRSP compared to G-LRSP.

We used instances with 25 customers and 30 customers in this experiment.
The LP relaxation of G-LRSP was solved using CPLEX 9.1 and AMPL
interface. The number of vehicles at each facility is set to number of
customers. This increases the number of variables, problem symmetry and
the LP solution times significantly. Because of this reason, the LP relaxation
values for some instances could not be obtained within the time limit. The
LP relaxation of SPP-LRSP (model (\ref{spp-obj}) through  (\ref{binary2})) was
obtained by using column generation algorithm described in Chapter 2. The algorithm
was coded in MINTO with CPLEX 9.1 as an LP solver.  Since in this experiment, we used
an earlier version of our column generation algorithm, we could not get
LP relaxation values for some of the instances within the time limit.
All instances in both experiment groups were run with a time limit of 8 CPU
hours.

The solution information for 25- and 30-customer instances are provided in Tables
\ref{comparesp-eb-gr1} and \ref{comparesp-eb-gr2}. Tables report LP relaxation
values and solution times for G-LRSP and SPP-LRSP for each instance.
Based on Table \ref{comparesp-eb-gr1} and Table \ref{comparesp-eb-gr2},
for all solved instances, the LP relaxations obtained from model SPP-LRSP are \%4.89 to
\%12.82 higher than the LP relaxations obtained from model G-LRSP.
In addition, the computation times for the SPP-LRSP are much better than the times
for G-LRSP in all solved instances. Furthermore, only 2 instances cannot be
solved within time limit using model SPP-LRSP, while the number of unsolved instances
is 8 if model G-LRSP is used.

\begin{table}[h]%[!hbp]
\centering
\caption{LP Relaxation Bounds of G-LRSP and SPP-LRSP: 25 Customer Instances}\label{comparesp-eb-gr1}
\begin {tabular}[h]{|c|c|c|c|c|c|}\hline
\mbox{} &  \multicolumn{2}{|c|}{ SPP-LRSP} & \multicolumn{2}{|c|}{G-LRSP} &  (a-b)/b*100\\\cline{2-5}
 Data File  & LP (a)    & CPU (s) &  LP (b) &   CPU (s) &  \% Gap   \\\hline
p01-f25-v1t1 &  3936.58 &   114.66  & 3705.06 & 5058.9  &   6.25    \\
p01-f25-v1t2 &  3855.91 &   346.18  & 3641.57 & 4719.67 &   5.89    \\
p01-f25-v2t1 &  3793.43 &   672.07  & 3610.19 & 4896.98 &   5.08    \\
p01-f25-v2t2 &  3724.91 &   4240.58 & 3551.19 & 5046.43 &   4.89    \\\hline
p03-f25-v1t1 &  4620.03 &   26.18   & 4267.66 & 8969.62 &   8.26    \\
p03-f25-v1t2 &  4538.07 &   256.77  & 4195.35 & 8729.64 &   8.17    \\
p03-f25-v2t1 &  4443.34 &   121.03  & 4179.83 & 10256.6 &   6.3 \\
p03-f25-v2t2 &  4351.62 &   791.77  & 4113.72 & 10934.8 &   5.78    \\\hline
p03-l25-v1t1 &  3928.73 &   34.83   & 3599.43 & 10906.3 &   9.15    \\
p03-l25-v1t2 &  3829.09 &   112.31  & 3527.39 & 10216.5 &   8.55    \\
p03-l25-v2t1 &  3812    &   78.6    & 3514.9  & 8722.87 &   8.45    \\
p03-l25-v2t2 &  3713.14 &   355.46  & 3445.92 & 8835.58 &   7.75    \\\hline
p07-f25-v1t1 &  3413.54 &   63.77   & 3151.23 & 15809.4 &   8.32    \\
p07-f25-v1t2 &  3341.53 &   394.39  & 3085.95 & 12698.5 &   8.28    \\
p07-f25-v2t1 &  3301.11 &   198.84  & 3069.45 & 13462.4 &   7.55    \\
p07-f25-v2t2 &  3198.43 &   677.1   & 3007.23 & 11832.1 &   6.36    \\\hline
p11-f25-v1t1 &  6615.99 &   0.42    & 5934.5  & 13097.1 &   11.48   \\
p11-f25-v1t2 &  6536.44 &   0.72    & 5865.22 & 14250.6 &   11.44   \\
p11-f25-v2t1 &  6378.77 &   0.73    & 5712.4  & 14334.6 &   11.67   \\
p11-f25-v2t2 &  6261.63 &   1.19    & 5650.04 & 14252.1 &   10.82   \\\hline
p11-l25-v1t1 &  7298.5  &   0.26    & 6469.19 & 9429.75 &   12.82   \\
p11-l25-v1t2 &  7187.6  &   0.33    & 6402.79 & 9334.2  &   12.26   \\
p11-l25-v2t1 &  6939.53 &   0.36    & 6188.64 & 11521.1 &   12.13   \\
p11-l25-v2t2 &  6904.43 &   0.44    & 6131.83 & 11561.7 &   12.6    \\\hline
\end{tabular}
\end{table}

\begin{table}[h]%[!hbp]
\centering
\caption{LP Relaxation Bounds of G-LRSP and SPP-LRSP: 30 Customer Instances}\label{comparesp-eb-gr2}
\begin {tabular}[h]{|c|c|c|c|c|c|}\hline
\mbox{} &  \multicolumn{2}{|c|}{SPP-LRSP} & \multicolumn{2}{|c|}{ G-LRSP} & (a-b)/b*100 \\\cline{2-5}
 Data File  & LP (a)    & CPU (s) & LP  (b) &  CPU (s) & \% Gap \\\hline
p01-f30-v1t1 &  4640.47 & 643.4     &   4287.79 &  17337.3  & 8.23    \\
p01-f30-v1t2 &  4561.62 & 8721.2    &   4219.16 &  17751.3  & 8.12    \\
p01-f30-v2t1 &  4462.66 & 4464.09   &   4169.85 &  21347.8  & 7.02    \\
p01-f30-v2t2 &  \mbox{-}& 29209.74  &   4106.6  &  18790    & \mbox{-} \\\hline
p03-f30-v1t1 &  5413.72 & 370.4     &   \mbox{-} & 28829.6  & \mbox{-} \\
p03-f30-v1t2 &  5330.33 & 6625.36   &   \mbox{-} & 28827.6  & \mbox{-} \\
p03-f30-v2t1 &  5209.41 & 2953.18   &   \mbox{-} & 28835    & \mbox{-} \\
p03-f30-v2t2 &  \mbox{-} & 29356.69 &   \mbox{-} & 28834.2  & \mbox{-} \\\hline
p03-l30-v1t1 &  4707.72 & 131.18    &   4317.71 &  20964.7  & 9.03  \\
p03-l30-v1t2 &  4615.22 & 2806.27   &   4240.26 &  19245.8  & 8.84  \\
p03-l30-v2t1 &  4566.1  & 521.17    &   4221.57 &  19880.2  & 8.16  \\
p03-l30-v2t2 &  4459.03 & 8129.62   &   4148.77 &  19549.6  & 7.48  \\\hline
p07-f30-v1t1 &  4153.16 & 705.12    &   3858.62 &  20347.3  & 7.63  \\
p07-f30-v1t2 &  4077.68 & 2625.43   &   3782.79 &  20362.7  & 7.8   \\
p07-f30-v2t1 &  3986.37 & 1870.87   &   3746.43 &  22084.3  & 6.4   \\
p07-f30-v2t2 &  3899.99 & 15385.17  &   3674.65 &  22243.4  & 6.13  \\\hline
p11-f30-v1t1 &  7232.02 & 0.59      &   6483.4  &  17474.6  & 11.55 \\
p11-f30-v1t2 &  7124.96 & 0.9       &   6410.9  &  18558.7  & 11.14 \\
p11-f30-v2t1 &  6965.07 & 1.22      &   6238.88 &  22240    & 11.64 \\
p11-f30-v2t2 &  6875.87 & 1.81      &   6172.69 &  23976.6  & 11.39 \\\hline
p11-l30-v1t1 &  8265.62 & 0.63      &   \mbox{-} & 28843    & \mbox{-} \\
p11-l30-v1t2 &  8161.91 & 1     &   \mbox{-} & 28823.6  & \mbox{-} \\
p11-l30-v2t1 &  8006.08 & 1.09      &   \mbox{-} & 28817.9  & \mbox{-} \\
p11-l30-v2t2 &  7798.55 & 2.53      &   \mbox{-} & 28816    & \mbox{-} \\\hline
\multicolumn{6}{l}{- Time limit is exceeded before LP solution is found.}
\end{tabular}
\end{table}

\subsection{Comparing  Alternative Set Partitioning Formulations}
\label{alternativeForm}
In this section, we compare the strength of the formulations constructed with
the valid inequalities described in Section \ref{StrongForm}.
We enumerated all of the feasible pairings and solved the instances using a
standard branch-and-bound algorithm without column generation.
Due to the large number of feasible pairings,
we had to conduct our experiments with smaller instances containing only 20 customers.
%In Table \ref{lpbounds-gr1}, for each instance, the first column lists the referenced
%\citet{Cordeau} instance, the second column lists the subset of customers,
%and the third and fourth columns list the vehicle capacity and time limit,
%respectively.

Let $SPP$ denote the base SPP-LRSP formulation that contains only constraints
(\ref{spp_con1}) and (\ref{spp_con2}).  We constructed $SPP_1$ by adding constraints
(\ref{relation-Z-T}), $SPP_2$ by adding constraints (\ref{relation}),
$SPP_3$ by adding constraint (\ref{lb-open-fac}), and $SPP_4$ by adding
both constraints (\ref{lb-open-fac}) and (\ref{relation}).
In Table~\ref{lpbounds-gr1}, for each instance, we report the total number of feasible pairings,
the integer optimal objective value, and the percentage gap of
the LP relaxation associated with each of the formulations.
In the table, a $-$ indicates that the branch-and-bound algorithm could not solve the
relaxation within 6 hours of CPU time.  For one instance, the set of feasible pairings
could not be generated, so we did not use it in the experiments.

We compared $SPP$ with $SPP_1$ and $SPP_2$ to evaluate the strength of constraints
(\ref{relation-Z-T}) and (\ref{relation}). Notice that adding inequalities (\ref{relation-Z-T})
in $SPP_1$ made the problem difficult to solve and yielded at most only a small improvement
in the LP relaxation bound for each solved instance.  Adding inequalities (\ref{relation}) in
$SPP_2$ improved the LP relaxation bound by a moderate amount.

Comparing $SPP$ with $SPP_3$ and $SPP_4$ evaluated the impact of adding constraint (\ref{lb-open-fac})
alone and in combination with constraints (\ref{relation}).
In all instances, adding inequality (\ref{lb-open-fac}) in $SPP_3$ lead to a dramatic
reduction in the gap as compared to adding (\ref{relation}) alone.
Together, the two inequalities (\ref{lb-open-fac}) and (\ref{relation}) in $SPP_4$ further reduced the
gap for all but 3 instances.  The average gap for $SPP_3$ is 3.54 while
the average gap for $SPP_4$ is 2.28.
These results support the theoretical results in Section \ref{StrongForm} and suggest that adding
both inequalities (\ref{lb-open-fac}) and (\ref{relation}) to the SPP-LRSP yields a stronger formulation,
which we denote as AUG-SPP-LRSP and used in Chapter 2 for the solution algorithms.
\begin{table}[!h]
\centering
\caption{Comparison of LP relaxation bounds with constraints (\ref{relation-Z-T}),
(\ref{relation}) and (\ref{lb-open-fac})} \label{lpbounds-gr1}
\begin {tabular}[!h]{|c|r|c|c|c|c|c|c|}
\hline
\mbox{} & \# of \ \  & \mbox{} & $SPP$  & $SPP_1$ & $SPP_2$  & $SPP_3$ & $SPP_4$\\
Data     & Pairing  & IP Obj  & \% Gap & \% Gap & \% Gap    & \% Gap  & \% Gap  \\\hline
p01-f20-v1t1    & 67558 & 4420  & 24.54 & 24.49 & 19.79 & 2.48  & 1.78  \\
p01-f20-v1t2    & 154903 & 4419 & 26.06 & \mbox{-}  & 22.08 & 3.69  & 3.52  \\
p01-f20-v2t1    & 147368 & 4378 & 26.17 & \mbox{-}  & 23.24 & 3.80  & 3.44  \\
p01-f20-v2t2    & 349885 & 4369 & 27.52 & \mbox{-}  & 25.05 & 4.86  & 4.83  \\\hline
p01-l20-v1t1    & 26266 & 4780  & 33.51 & 33.38 &   27.98   & 5.37  & 2.16  \\
p01-l20-v1t2    & 60211 & 4599  & 32.77 & 32.72 &   28.08   & 3.49  & 0.58  \\
p01-l20-v2t1    & 51538 & 4605  & 33.64 & 33.51 &   29.91   & 4.30  & 1.57  \\
p01-l20-v2t2    & 121027 & 4481 & 33.24 & 33.23 &   30.81   & 2.97  & 1.54  \\\hline
p03-f20-v1t1    & 52619 & 4468  & 15.64 & 15.63 &   14.08   & 0.72  & 0.43  \\
p03-f20-v1t2    & 121933 & 4461 & 17.25 & \mbox{-}  &   16.12   & 2.32  & 2.13  \\
p03-f20-v2t1    & 107172 & 4421 & 18.04 & \mbox{-}  &   17.42   & 2.88  & 2.75  \\
p03-f20-v2t2    & 254863 & 4421 & 19.78 & \mbox{-}  &   19.12   & 4.58  & 4.44  \\\hline
p03-m20-v1t1    & 37751 & 4543  & 18.6  & 18.42 &   15.83   &   3.20    &   3.02    \\
p03-m20-v1t2    & 89465 & 4442  & 18.13 & \mbox{-}  &   16.02   & 2.35  & 2.35  \\
p03-m20-v2t1    & 78580 & 4383  & 19.15 & \mbox{-}  &   17.59   & 3.19  & 3.19  \\
p03-m20-v2t2    & 193509 & 4376 & 20.13 & \mbox{-}  &   19.2    & 4.12  & 4.12  \\\hline
p03-l20-v1t1    & 32801 & 4617  & 31.27 & 30.84 &   26  & 4.74  & 1.72  \\
p03-l20-v1t2    & 79215 & 4467  & 30.19 & 30.13 &   27.08   & 2.77  & 1.30  \\
p03-l20-v2t1    & 51052 & 4404  & 29.84 & 29.84 &   28.35   & 1.45  & 1.04  \\
p03-l20-v2t2    & 133853 & 4399 & 31.76 & \mbox{-}  &   30.45   & 3.55  & 3.11  \\\hline
p07-f20-v1t1    & 68267 & 4426  & 38.22 & 37.96 &   34.78   & 2.81  & 1.13  \\
p07-f20-v1t2    & 166551 & 4421 & 39.33 & \mbox{-}  &   37.18   & 3.60  & 2.57  \\
p07-f20-v2t1    & 133873 & 4380 & 39.45 & \mbox{-}  &   38.55   & 3.25  & 2.58  \\\hline
p07-s20-v1t1    & 27623 & 4840  & 31.69 & 31.47 &   25.75   & 4.73  & 1.01  \\
p07-s20-v1t2    & 61753 & 4801  & 32.79 & 32.72 &   27.97   & 5.61  & 2.91  \\
p07-s20-v2t1    & 54964 & 4751  & 32.8  & 32.48 &   29.09   & 5.34  & 2.32  \\
p07-s20-v2t2    & 127679 & 4543 & 31.15 & \mbox{-}  &   28.1    & 2.43  & 0.68  \\\hline
p07-t20-v1t1    & 11267 & 4758  & 29.49 & 29.47 &   26.03   & 3.62  & 1.05  \\
p07-t20-v1t2    & 28315 & 4670  & 30.44 & 30.06 &   28.58   & 3.77  & 2.08  \\
p07-t20-v2t1    & 19283 & 4688  & 32.39 & 31.83 &   30.98   & 5.57  & 4.14  \\
p07-t20-v2t2    & 53689 & 4460  & 30.47 & 30.28 &   29.2    & 2.11  & 1.33  \\\hline
Avg. &  &  & 28.24   & 29.91 &  25.50  &  3.54 &   2.28   \\\hline
\end{tabular}
\end{table}

\begin{comment}
\begin{table}[h]
\centering
\caption{Comparison of LP relaxation bounds with constraints (\ref{relation-Z-T}),
(\ref{relation}) and (\ref{lb-open-fac})} \label{lpbounds-gr1}
\begin {tabular}[t]{|c|c|c|c|r|c|c|c|c|c|c|}
\hline
        &          & \mbox{} & \mbox{} & \# of \ \  & \mbox{} & $SPP$  & $SPP_1$ & $SPP_2$  & $SPP_3$ & $SPP_4$\\
Ref     & Cust     & $L^V$   & $L^T$   & Pairing  & IP Obj  & \% Gap & \% Gap & \% Gap    & \% Gap  & \% Gap  \\\hline
p01 & 1 - 20    & 50 & 140 & 67558  & 4420  & 24.54 & 24.49 & 19.79 & 2.48  & 1.78  \\
    &           & 50 & 160 & 154903 & 4419  & 26.06 & \mbox{-}  & 22.08 & 3.69  &   3.52    \\
    &           & 70 &  140 & 147368    & 4378  & 26.17 & \mbox{-}  & 23.24 & 3.80  &   3.44    \\
    &           & 70 &  160 & 349885    & 4369  & 27.52 & \mbox{-}  & 25.05 & 4.86  &   4.83    \\\hline
p01 & 31 - 50 &  50 &   140 & 26266 & 4780  & 33.51 & 33.38 &   27.98   & 5.37  &   2.16    \\
    & & 50 & 160 & 60211    & 4599  & 32.77 & 32.72 &   28.08   & 3.49  &   0.58    \\
    & & 70 &    140 & 51538 & 4605  & 33.64 & 33.51 &   29.91   & 4.30  &   1.57    \\
    & & 70 &    160 & 121027    & 4481  & 33.24 & 33.23 &   30.81   & 2.97  &   1.54    \\\hline
p03 & 1 - 20 & 60 & 140 & 52619 & 4468  & 15.64 & 15.63 &   14.08   & 0.72  &   0.43    \\
    &   & 60 &  160 & 121933    & 4461  & 17.25 & \mbox{-}  & 16.12 & 2.32  &   2.13    \\
    &       & 80 &  140 & 107172    & 4421  & 18.04 & \mbox{-}  & 17.42 & 2.88  &   2.75    \\
    &       & 80    & 160   & 254863 & 4421 & 19.78 & \mbox{-}  & 19.12 & 4.58  &   4.44    \\\hline
p03 & 28 - 47   & 60    & 140   & 37751 & 4543  & 18.6  & 18.42 &   15.83   & 3.20  &   3.02    \\
    &   & 60    & 160   & 89465 & 4442  & 18.13 & \mbox{-}  & 16.02 & 2.35  &   2.35    \\
    &       & 80    & 140   & 78580 & 4383  & 19.15 & \mbox{-}  & 17.59 & 3.19  &   3.19    \\
    &       & 80    & 160   & 193509 & 4376 & 20.13 & \mbox{-}  & 19.2  & 4.12  &   4.12    \\\hline
p03 & 56 - 75   & 60    & 140   & 32801 & 4617  & 31.27 & 30.84 &   26  & 4.74  &   1.72    \\
    & & 60  & 160   & 79215 & 4467  & 30.19 & 30.13 &   27.08   & 2.77  &   1.30    \\
    & & 80  & 140   & 51052 & 4404  & 29.84 & 29.84 &   28.35   & 1.45  &   1.04    \\
    & & 80  & 160   & 133853 & 4399 & 31.76 & \mbox{-}  & 30.45 & 3.55  &   3.11    \\\hline
p07 & 1 - 20 & 50   & 140   & 68267 & 4426  & 38.22 & 37.96 &   34.78   & 2.81  &   1.13    \\
    &   & 50    & 160   & 166551 & 4421 & 39.33 & \mbox{-}  & 37.18 & 3.60  &   2.57    \\
    & & 70  & 140   & 133873 & 4380 & 39.45 & \mbox{-}  & 38.55 & 3.25  &   2.58    \\\hline
p07 & 21 - 40   & 50    & 140   & 27623 & 4840  & 31.69 & 31.47 &   25.75   & 4.73  &   1.01    \\
    & & 50  & 160   & 61753 & 4801  & 32.79 & 32.72 &   27.97   & 5.61  &   2.91    \\
    & & 70  & 140   & 54964 & 4751  & 32.8  & 32.48 &   29.09   & 5.34  &   2.32    \\
    & & 70  & 160   & 127679 & 4543 & 31.15 & \mbox{-}  & 28.1  & 2.43  &   0.68    \\\hline
p07 & 31 - 50   & 50    & 140   & 11267 & 4758  & 29.49 & 29.47 &   26.03   & 3.62  &   1.05    \\
    & & 50  & 160   & 28315 & 4670  & 30.44 & 30.06 &   28.58   & 3.77  &   2.08    \\
    & & 70  & 140   & 19283 & 4688  & 32.39 & 31.83 &   30.98   & 5.57  &   4.14    \\
    & & 70  & 160   & 53689 & 4460  & 30.47 & 30.28 &   29.2    & 2.11  &   1.33    \\\hline
Avg. &   &  &  &  &  & 28.24   & 29.91 &  25.50  &  3.54 &   2.28   \\\hline
\end{tabular}
\end{table}

\begin{table}[h]
\centering
\caption{Comparison of LP relaxation bounds with constraints (\ref{relation-Z-T}) and (\ref{relation})}\label{lpbounds-gr1}
\begin {tabular}[t]{|c|c|c|c|c|c|c|c|c|c|c|}
\hline
Data &  \# of  &  IP & \multicolumn{2}{|c|}{$SPP_B$}& \multicolumn{3}{|c|}{$SPP_1$} & \multicolumn{3}{|c|}{$SPP_2$}\\ \cline {4-11}
\mbox{} &  Col. &  \mbox{} &  LP &  Gap  & LP & Impr. & Gap &  LP & Impr. & Gap \\
\mbox{} &  \mbox{} &  \mbox{} &  \mbox{} & \mbox{} & \mbox{} & (\%) & \mbox{} & \mbox{} &  (\%) & \mbox{} \\\hline

p01-f20-v1t1 &  67558 & 4420    &   3335.23 &   24.54   &   3337.72 &   0.07    &   24.49   &   3545.31 & 6.3   & 19.79 \\
p01-f20-v1t2 &  154903 &    4419    &   3267.48 &   26.06   &   \mbox{*}    &   \mbox{*}    &   \mbox{*} &  3443.35 & 5.38  & 22.08 \\
p01-f20-v2t1 &  147368 &    4378    &   3232.45 &   26.17   &   \mbox{*}    &   \mbox{*}    &   \mbox{*} &  3360.59 & 3.96  &   23.24   \\
p01-f20-v2t2 &  349885 &    4369    &   3166.44 &   27.52   &   \mbox{*}    &   \mbox{*}    &   \mbox{*} &  3274.78 & 3.42  &   25.05   \\
p01-l20-v1t1 &  26266 & 4780    &   3178.37 &   33.51   &   3184.38 &   0.19    &   33.38   &   3442.39 &   8.31    & 27.98 \\
p01-l20-v1t2 &  60211 & 4599    &   3091.84 &   32.77   &   3094.28 &   0.08    &   32.72   &   3307.59 &   6.98    & 28.08 \\
p01-l20-v2t1 &  51538 & 4605    &   3056.11 &   33.64   &   3061.8  &   0.19    &   33.51   &   3227.44 &   5.61    & 29.91 \\
p01-l20-v2t2 &  121027 &    4481    &   2991.67 &   33.24   &   2991.89 &   0.01    &   33.23   &   3100.46 &   3.64    &   30.81   \\
p03-f20-v1t1 &  52619 & 4468    &   3769.21 &   15.64   &   3769.55 &   0.01    &   15.63   &   3838.82 &   1.85    & 14.08 \\
p03-f20-v1t2 &  121933 &    4461    &   3691.56 &   17.25   &   \mbox{*}    &   \mbox{*}    &   \mbox{*}    & 3741.88   &   1.36    &   16.12   \\
p03-f20-v2t1 &  107172 &    4421    &   3623.24 &   18.04   &   \mbox{*}    &   \mbox{*}    &   \mbox{*}    & 3650.7    &   0.76    &   17.42   \\
p03-f20-v2t2 &  254863 &    4421    &   3546.37 &   19.78   &   \mbox{*}    &   \mbox{*}    &   \mbox{*}    & 3575.71   &   0.83    &   19.12   \\
p03-m20-v1t1 &  37751 & 4543    &   3697.98 &   18.6    &   3706.08 &   0.22    &   18.42   &   3823.66 &   3.4 & 15.83 \\
p03-m20-v1t2 &  89465 & 4442    &   3636.45 &   18.13   &   \mbox{*}    &   \mbox{*}    &   \mbox{*}    &   3730.25 &   2.58    &   16.02   \\
p03-m20-v2t1 &  78580 & 4383    &   3543.66 &   19.15   &   \mbox{*}    &   \mbox{*}    &   \mbox{*}    &   3612.2  &   1.93    &   17.59   \\
p03-m20-v2t2    &   193509  &   4376    &   3495.18 &   20.13   &   \mbox{*}    &   \mbox{*}    &   \mbox{*} &  3535.99 &   1.17    &   19.2    \\
p03-l20-v1t1 &  32801 & 4617    &   3173.4  &   31.27   &   3192.92 &   0.62    &   30.84   &   3416.8  &   7.67    & 26    \\
p03-l20-v1t2 &  79215 & 4467    &   3118.36 &   30.19   &   3121.16 &   0.09    &   30.13   &   3257.28 &   4.45    & 27.08 \\
p03-l20-v2t1 &  51052 & 4404    &   3089.98 &   29.84   &   3090.05 &   0   &   29.84   &   3155.67 &   2.13    & 28.35 \\
p03-l20-v2t2 &  133853 &    4399    &   3001.81 &   31.76   &   \mbox{*}    &   \mbox{*}    &   \mbox{*}    & 3059.34   &   1.92    &   30.45   \\
p07-f20-v1t1 &  68267 & 4426    &   2734.17 &   38.22   &   2745.73 &   0.42    &   37.96   &   2886.51 &   5.57    & 34.78 \\
p07-f20-v1t2 &  166551 &    4421    &   2682.31 &   39.33   &   \mbox{*}    &   \mbox{*}    &   \mbox{*}    & 2777.34   &   3.54    &   37.18   \\
p07-f20-v2t1 &  133873 &    4380    &   2652.23 &   39.45   &   \mbox{*}    &   \mbox{*}    &   \mbox{*}    & 2691.56   &   1.48    &   38.55   \\
p07-f20-v2t2 &  \mbox{**}   &   4380    &   \mbox{-}    &   \mbox{-}    &   \mbox{-}    &   \mbox{-}    & \mbox{-}  &   \mbox{-}    &   \mbox{-}    &   \mbox{-}    \\
p07-s20-v1t1 &  27623 & 4840    &   3306.25 &   31.69   &   3316.8  &   0.32    &   31.47   &   3593.69 &   8.69    & 25.75 \\
p07-s20-v1t2 &  61753 & 4801    &   3226.85 &   32.79   &   3229.96 &   0.1 &   32.72   &   3458.14 &   7.17    & 27.97 \\
p07-s20-v2t1 &  54964 & 4751    &   3192.57 &   32.8    &   3208.11 &   0.49    &   32.48   &   3369.09 &   5.53    & 29.09 \\
p07-s20-v2t2 &  127679 &    4543    &   3127.67 &   31.15   &   \mbox{*}    &   \mbox{*}    &   \mbox{*}    & 3266.19   &   4.43    &   28.1    \\
p07-t20-v1t1 &  11267 & 4758    &   3354.79 &   29.49   &   3355.91 &   0.03    &   29.47   &   3519.27 &   4.9 & 26.03 \\
p07-t20-v1t2 &  28315 & 4670    &   3248.33 &   30.44   &   3266.37 &   0.56    &   30.06   &   3335.52 &   2.68    & 28.58 \\
p07-t20-v2t1 &  19283 & 4688    &   3169.33 &   32.39   &   3195.69 &   0.83    &   31.83   &   3235.66 &   2.09    & 30.98 \\
p07-t20-v2t2 &  53689 & 4460    &   3101.05 &   30.47   &   3109.41 &   0.27    &   30.28   &   3157.5  &   1.82    & 29.2  \\\hline
\multicolumn {11}{l}{* not be able to solve LP in 6 hours of CPU time} \\
\multicolumn {11}{l}{** not be able to generate all columns in 10 days of CPU time} \\
\end{tabular}
\end{table}

\begin{table}[h]
\centering
\caption{Comparison of LP relaxation bounds with constraints (\ref{lb-open-fac}) and (\ref{relation})}\label{lpbounds-gr2}
\begin {tabular}[t]{|c|c|c|c|c|c|c|c|c|c|}
\hline
{\bf Instance} & {\bf IP} & \multicolumn{2}{|c|}{{\bf $SPP_B$}} & \multicolumn{3}{|c|}{{\bf $SPP_3$}} & \multicolumn{3}{|c|}{{\bf $SPP_4$}}\\ \cline {3-10}
\mbox{}  &  \mbox{} &  {\bf LP} & {\bf LP Rlx.} &{\bf LP} & {\bf Impr.} & {\bf LP Rlx.} & {\bf LP} & {\bf Impr.} & {\bf LP Rlx.} \\
\mbox{} &   \mbox{} &  \mbox{} & {\bf gap(\%)} & \mbox{} & {\bf (\%)} & {\bf gap(\%)} & \mbox{} & {\bf (\%)} & {\bf gap(\%)} \\\hline
p01-f20-v1t1    &   4420    &   3335.23 &   24.54   &   4310.45 &   29.24   &   2.48    &   4341.31 &   30.17   &   1.78    \\
p01-f20-v1t2    &   4419    &   3267.48 &   26.06   &   4256.15 &   30.26   &   3.69    &   4263.26 &   30.48   &   3.52    \\
p01-f20-v2t1    &   4378    &   3232.45 &   26.17   &   4211.65 &   30.29   &   3.8 &   4227.25 &   30.78   &   3.44    \\
p01-f20-v2t2    &   4369    &   3166.44 &   27.52   &   4156.57 &   31.27   &   4.86    &   4157.86 &   31.31   &   4.83    \\
p01-l20-v1t1    &   4780    &   3178.37 &   33.51   &   4523.2  &   42.31   &   5.37    &   4676.72 &   47.14   &   2.16    \\
p01-l20-v1t2    &   4599    &   3091.84 &   32.77   &   4438.37 &   43.55   &   3.49    &   4572.51 &   47.89   &   0.58    \\
p01-l20-v2t1    &   4605    &   3056.11 &   33.64   &   4406.91 &   44.2    &   4.3 &   4532.65 &   48.31   &   1.57    \\
p01-l20-v2t2    &   4481    &   2991.67 &   33.24   &   4347.86 &   45.33   &   2.97    &   4411.9  &   47.47   &   1.54    \\
p03-f20-v1t1    &   4468    &   3769.21 &   15.64   &   4436.01 &   17.69   &   0.72    &   4448.7  &   18.03   &   0.43    \\
p03-f20-v1t2    &   4461    &   3691.56 &   17.25   &   4357.69 &   18.04   &   2.32    &   4365.97 &   18.27   &   2.13    \\
p03-f20-v2t1    &   4421    &   3623.24 &   18.04   &   4293.76 &   18.51   &   2.88    &   4299.58 &   18.67   &   2.75    \\
p03-f20-v2t2    &   4421    &   3546.37 &   19.78   &   4218.54 &   18.95   &   4.58    &   4224.68 &   19.13   &   4.44    \\
p03-m20-v1t1    &   4543    &   3697.98 &   18.6    &   4397.78 &   18.92   &   3.2 &   4405.89 &   19.14   &   3.02    \\
p03-m20-v1t2    &   4442    &   3636.45 &   18.13   &   4337.65 &   19.28   &   2.35    &   4337.77 &   19.29   &   2.35    \\
p03-m20-v2t1    &   4383    &   3543.66 &   19.15   &   4243.33 &   19.74   &   3.19    &   4243.33 &   19.74   &   3.19    \\
p03-m20-v2t2    &   4376    &   3495.18 &   20.13   &   4195.8  &   20.05   &   4.12    &   4195.8  &   20.05   &   4.12    \\
p03-l20-v1t1    &   4617    &   3173.4  &   31.27   &   4398.3  &   38.6    &   4.74    &   4537.44 &   42.98   &   1.72    \\
p03-l20-v1t2    &   4467    &   3118.36 &   30.19   &   4343.05 &   39.27   &   2.77    &   4409.02 &   41.39   &   1.3 \\
p03-l20-v2t1    &   4404    &   3089.98 &   29.84   &   4340.22 &   40.46   &   1.45    &   4358.23 &   41.04   &   1.04    \\
p03-l20-v2t2    &   4399    &   3001.81 &   31.76   &   4242.69 &   41.34   &   3.55    &   4262    &   41.98   &   3.11    \\
p07-f20-v1t1    &   4426    &   2734.17 &   38.22   &   4301.54 &   57.33   &   2.81    &   4376.1  &   60.05   &   1.13    \\
p07-f20-v1t2    &   4421    &   2682.31 &   39.33   &   4261.95 &   58.89   &   3.6 &   4307.48 &   60.59   &   2.57    \\
p07-f20-v2t1    &   4380    &   2652.23 &   39.45   &   4237.74 &   59.78   &   3.25    &   4267    &   60.88   &   2.58    \\
p07-f20-v2t2    &   4380    &   \mbox{-}    & \mbox{-}  &   \mbox{-}    &   \mbox{-}    &   \mbox{-}    &   \mbox{-}    &   \mbox{-}    &   \mbox{-}    \\
p07-s20-v1t1    &   4840    &   3306.25 &   31.69   &   4611.17 &   39.47   &   4.73    &   4790.9  &   44.9    &   1.01    \\
p07-s20-v1t2    &   4801    &   3226.85 &   32.79   &   4531.77 &   40.44   &   5.61    &   4661.13 &   44.45   &   2.91    \\
p07-s20-v2t1    &   4751    &   3192.57 &   32.8    &   4497.49 &   40.87   &   5.34    &   4641    &   45.37   &   2.32    \\
p07-s20-v2t2    &   4543    &   3127.67 &   31.15   &   4432.59 &   41.72   &   2.43    &   4512.13 &   44.26   &   0.68    \\
p07-t20-v1t1    &   4758    &   3354.79 &   29.49   &   4585.89 &   36.7    &   3.62    &   4708.04 &   40.34   &   1.05    \\
p07-t20-v1t2    &   4670    &   3248.33 &   30.44   &   4493.77 &   38.34   &   3.77    &   4573.04 &   40.78   &   2.08    \\
p07-t20-v2t1    &   4688    &   3169.33 &   32.39   &   4427.09 &   39.69   &   5.57    &   4493.88 &   41.79   &   4.14    \\
p07-t20-v2t2    &   4460    &   3101.05 &   30.47   &   4365.97 &   40.79   &   2.11    &   4400.83 &   41.91   &   1.33    \\
\hline
\multicolumn {10}{l}{** not be able to generate all columns in 10 days of CPU time} \\
\end{tabular}
\end{table}
\end{comment}


\section{Computational Experiments For Chapter 2}
The goal of the computational experiments in this section
is to demonstrate the effect of the algorithms for the components of a branch and price
algorithm  (described in Chapter 2) and to reason the extensions done to the branch and price
algorithms to improve the overall performance. We evaluate the different pricing algorithms
used in column generation process, effect of updating column generation with heuristic
pricing algorithms, the strength of the primal heuristics to update the upper bound, the
effect of the of the initial set of columns to the overall branch and price algorithm,
\zaupdate{and the effect of branching rule selection procedure}.
Finally, the performance of the branch and price algorithms is demonstrated for
25- and 40-customer instances.

\subsection{Performance of Exact Pricing Algorithms: 1p-ESPPRC and 2p-ESPPRC}
In this section, we validate the claimed performance changes with different exact
pricing problem solution algorithms in Section \ref{solv-pri-exact}
by empirically demonstrating the performance change.

We solved the LP relaxation of AUG-SPP-LRSP using column generation algorithm
(referred as COLGEN) with one of the pricing algorithms: one phase pricing problem
solution algorithm (referred as 1p-ESPPRC), one phase pricing problem solution algorithm
with new domination rule (presented in Section (\ref{1pespprc})) (referred as 1p-ESPPRC-NwD),
and two phase pricing problem solution algorithm (referred as 2p-ESPPRC). In the experiments,
at each iteration of COLGEN, at most one column for each facility (one with most negative
reduced cost) was added to the RMP. In addition, we tested the effect of adding multiple
columns instead of at most one at each iteration of COLGEN. We added at most 40 columns
with most negative reduced cost for each facility at each iteration of COLGEN with 2p-ESPPRC.
%We ran the experiments with a time limit of 3 CPU hours.

In Table \ref{compexp-comp-pri1}, for each instance and for each design of COLGEN,
we report the total CPU time (in seconds) to solve the LP relaxation of RMP (at root node)
and average time spent at each COLGEN iteration (total time / number of column generation
iterations). For 2p-ESPPRC and 2p-ESPPRC with 40
columns, we also report number of COLGEN iterations. In Table \ref{compexp-comp-pri1},
we calculate the fraction of total time (for 1p-ESPPRC-NwD, 2p-ESPPRC and 2p-ESPPRC with 40
columns) with respect to total time for 1p-ESPPRC.

Results on Table \ref{compexp-comp-pri1} show that new domination rule significantly
improves the performance of 1p-ESPPRC. On average, total COLGEN time for 1p-ESPPRC-NwD is
\%64 of the total time in 1p-ESPPRC.  Average time per COLGEN iteration improves
about \%20 to \%50 with the addition of new domination rule to 1p-ESPPRC.

As seen from Table \ref{compexp-comp-pri1}, 2p-ESPPRC reduces the total time \%92 on
average. And average time per COLGEN iteration reduces significantly for every instance.
Further improvement is obtained by adding 40 columns at each iteration
in 2p-ESPPRC. Adding multiple columns at each iteration instead of one column reduces number of
COLGEN iterations in every instance. On average, total time for 2p-ESPPRC with multiple columns
is \%49 better than the time with single column 2p-ESPPRC.

Based on these experiments, 2p-ESPPRC algorithm performs significantly better than 1p-ESPPRC
and 1p-ESPPRC-NwD to solve exact pricing problem. In addition, benefit of adding multiple columns
instead of single best column at each iteration  empirically demonstrated with this experiment.
\begin{table}[h]
\centering
\caption{Performance of 1p-ESPPRC, 1p-ESPPRC-NwD, 2p-ESSPRC}\label{compexp-comp-pri1}
\begin {tabular}[t]{|c|c|c|c|c|c|}
\hline
\mbox{} & \multicolumn{2}{|c|}{1p-ESPPRC} & \multicolumn{3}{|c|}{1p-ESPPRC-NwD} \\\cline {2-6}
\mbox{} &  Tot.time  & Time per    & Tot.time  & (b/a)   & Time per    \\
Data File & (s) (a)& iter        & (s) (b) &         & iter \\\hline
p01-f25-v2t2  & 13467.95&   134.68  &   6515.66 &   0.48    &   69.32 \\
p01-l25-v2t2  & 2127.41 &   31.75   &   1601.02 &   0.75    &   24.63 \\
p03-f25-v2t2  & 3273.45 &   32.73   &   2054.31 &   0.63    &   23.08  \\
p03-m25-v2t2  & 12680.8 &   162.57  &   6380.45 &   0.5 &   76.87 \\
p03-l25-v2t2  & 1417.49 &   19.96   &   962.73  &   0.68    &   13.56 \\
p07-f25-v2t2  & 2231.63 &   31.88   &   1664.52 &   0.75    &   24.48 \\
p07-s25-v2t2  & 305.07  &   4.49    &   220.66  &   0.72    &   3.56 \\
p07-t25-v2t2  & 4036.72 &   61.16   &   2492.33 &   0.62    &   39.56 \\\hline
Average       & 4942.57 &   59.9    &   2736.46 &   0.64    &   34.38 \\\hline
\multicolumn{6}{c}{\mbox{}}\\
\multicolumn{6}{c}{\mbox{}}\\
\end{tabular}
\\
\begin {tabular}[t]{|c|c|c|c|c|c|c|c|c|}
\hline
\mbox{} & \multicolumn{4}{|c|}{2p-ESPPRC} & \multicolumn{4}{|c|}{2p-ESPPRC with 40 col.per iter.} \\\cline {2-9}
\mbox{} &  Tot.time   & (c/a)  & \# of & Time per    & Tot.time  & (d/c)  & \# of & Time per \\
Data File & (s) (c) &        & iter  & iter        & (s) (d) &        & iter  & iter \\\hline
p01-f25-v2t2    & 515.59 &  0.04    & 94    & 5.49  & 143.92 & 0.28 & 52    & 2.77  \\
p01-l25-v2t2    & 195.69 &  0.09    & 70    & 2.8   & 95.13  & 0.49 & 29    & 3.28  \\
p03-f25-v2t2    & 142.52 &  0.04    & 81    & 1.76  & 70.43  & 0.49 & 35    & 2.01  \\
p03-m25-v2t2    & 321.98 &  0.03    & 78    & 4.13  & 173.35 & 0.54 & 38    & 4.56  \\
p03-l25-v2t2    & 114.12 &  0.08    & 75    & 1.52  & 76.26  & 0.67 & 33    & 2.31  \\
p07-f25-v2t2    & 293.53 &  0.13    & 64    & 4.59  & 140.89 & 0.48 & 27    & 5.22  \\
p07-s25-v2t2    & 51.44  &  0.17    & 59    & 0.87  & 27.7   & 0.54 & 31    & 0.89  \\
p07-t25-v2t2    & 223.61 &  0.06    & 68    & 3.29  & 124.55 & 0.56 & 29    & 4.29  \\\hline
Average         & 232.31 &  0.08    & 73.63 & 3.05  & 106.53 & 0.51 & 34.25 & 3.17  \\\hline
\end{tabular}
\end{table}

\subsection{Effect of Adding Heuristic Pricing Algorithms: ESPPRC-LL and ESPPRC-CS}
In this section, we empirically demonstrate the performance
improvement that can be obtained by addition of heuristic pricing algorithms
to overall column generation algorithm. In Section \ref{heur-pri-problems}, we describe
two heuristic versions of pricing algorithm: ESPPRC with label limit (ESPPRC-LL) and
ESPPRC for a subset of customers (ESPPRC-CS).

We applied branch and price algorithm using different versions of column generation algorithm
to solve the instances. We designed four experiments. In the first experiment, at each node
of the branch and price algorithm,
we used only exact pricing, 2p-ESPPRC in COLGEN. In the second experiment, at each node, we ran
heuristic pricing algorithm 2p-ESPPRC-LL with label limits 50, 100 and 200 before the exact
pricing algorithm. At each node we solved 2p-ESPPRC-LL(50) until no column
is found. Then we incremented the label limit to 100 and applied 2p-ESPPRC-LL(100). When
2p-ESPPRC-LL(100) terminated with empty column list, we started applying  2p-ESPPRC-LL(200).
Finally, we applied 2p-ESPPRC until the column generation terminates.
In the third experiment, we applied 2p-ESPPRC-CS with subset size 12. 2p-ESPPRC-CS(12)
is ran at most 12 iterations before 2p-ESPPRC was called. In the last experiment,
we applied both heuristic pricing algorithms before exact pricing algorithm. We applied
2p-ESPPRC-CS(12), then 2p-ESPPRC-LL with label limits 50, 100 and 200, and finally 2p-ESPPRC.


In Table \ref{compexp-comp-pri2}, for each instance and for each design of branch and
price algorithm, we report average column generation time (total COLGEN time / \# of evaluated
nodes), \# of evaluated nodes and total algorithm time (in CPU seconds). In Table
\ref{compexp-comp-pri2}, we calculate the fraction of total time in each version of
the branch and price algorithm with respect to total time of the algorithm
with only 2p-ESPPRC (without any heuristic pricing algorithms).

Results on Table \ref{compexp-comp-pri2} show that applying 2p-ESPPRC-LL before exact pricing
algorithm improves the performance of the overall algorithm, i.e., total algorithm time improves
\%20 on average. In all instances, total solution time and average COLGEN time is same or
lower in branch and price algorithm with 2p-ESPPRC-LL.

The algorithm with  2p-ESPPRC-CS improves the solution time only \%2 on average. In half of
the test instances, solution time with ESPPRC-CS is better than the algorithm with only
exact pricing. The algorithm using heuristic pricing algorithms 2p-ESPPRC-CS and 2p-ESPPRC-LL
in combination is about \%17 better on average. This result supports that \%20 effect of
2p-ESPPRC-LL in the second experiment. Therefore, we can conclude that in the fourth experiment
the performance improvement was obtained by application of 2p-ESPPRC-LL, and
2p-ESPPRC-CS seems to have a negative effect to the performance. Even though the contribution
of 2p-ESPPRC-CS seems negligible, the effect of heuristic pricing 2p-ESPPRC-CS is more significant
with larger instances. We have implemented another set of experiments to show this effect. We
describe it next.
\begin{table}[h]
\centering
\caption{Performance improvement with 2p-ESPPRC-LL and 2p-ESSPRC-CS}\label{compexp-comp-pri2}
\begin {tabular}[h]{|c|c|c|c|c|c|c|c|}
\hline
\mbox{} & \multicolumn{3}{|c|}{BP with 2p-ESPPRC} & \multicolumn{4}{|c|}{BP with 2p-ESPPRC-LL + 2p-ESPPRC} \\\cline {2-8}
\mbox{} &  Avg. CG   & \# of    & Tot.time  &  Avg. CG   & \# of    & Tot.time  & (b/a) \\
Data File & time (s) & BB Nd.   & (s) (a) &  time (s)  & BB Nd.   & (s) (b) &  \mbox{} \\\hline
p01-f25-v2t2    & 10.86 & 151   & 1670.57 & 7.79    & 119   & 954.41    & 0.57  \\
p01-l25-v2t2    & 14.25 & 37    & 533.13  & 10.48   & 39    & 414.51    & 0.78  \\
p03-f25-v2t2    & 6.41  & 49    & 322.28  & 5.92    & 52    & 316.88    & 0.98  \\
p03-m25-v2t2    & 92.85 & 9 & 838.71  & 57.89   & 9 & 523.75    & 0.62  \\
p03-l25-v2t2    & 13.76 & 41    & 568.25  & 12.11   & 37    & 452.17    & 0.8   \\
p07-f25-v2t2    & 8.84  & 135   & 1214.31 & 4.39    & 197   & 907.33    & 0.75  \\
p07-s25-v2t2    & 2.05  & 102   & 217.42  & 2.1         & 100   & 217.97    & 1 \\
p07-t25-v2t2    & 9.08  & 87    & 816.85  & 7.92    & 91    & 754.9     & 0.92  \\\hline
Average     & 19.76 & 76.38 & 772.69  & 13.58   & 80.5  & 567.74    & 0.8   \\\hline
\multicolumn{6}{c}{\mbox{}}\\
\multicolumn{6}{c}{\mbox{}}\\
\end{tabular}
\\
\begin {tabular}[h]{|c|c|c|c|c|c|c|c|c|}
\hline
\mbox{} & \multicolumn{4}{|c|}{BP with 2p-ESPPRC-CS + } & \multicolumn{4}{|c|}{BP with 2p-ESPPRC-LL + }\\ %\cline {2-9}
\mbox{} & \multicolumn{4}{|c|}{2p-ESPPRC} & \multicolumn{4}{|c|}{2p-ESPPRC-CS + 2p-ESPPRC}\\\cline {2-9}
\mbox{} &  Avg. CG   & \# of    & Tot.time  &  (c/a) &   Avg. CG   & \# of    & Tot.time  & (d/a) \\
Data File & time (s) & BB Nd.   & (s) (c) &        &   time (s)  & BB Nd.   & (s) (d) &  \mbox{} \\\hline
p01-f25-v2t2    & 8.68  & 101 & 897.41 & 0.54   & 11.78 & 71    & 849.63 & 0.51 \\
p01-l25-v2t2    & 14.63 & 37  & 546.99 & 1.03   & 12.12 & 39    & 478.4  & 0.9  \\
p03-f25-v2t2    & 7.1   & 50  & 365.63 & 1.13   & 6.4   & 44    & 289.4  & 0.9  \\
p03-m25-v2t2    & 63.71 & 9   & 576.48 & 0.69   & 49.05 & 9 & 444.39 & 0.53 \\
p03-l25-v2t2    & 14.06 & 39  & 552.7   & 0.97  & 12.18 & 37    & 454.43 & 0.8  \\
p07-f25-v2t2    & 4.71  & 219 & 1076.36 & 0.89  & 3.76  & 235   & 938.72 & 0.77 \\
p07-s25-v2t2    & 2.3   & 94 & 223.8    & 1.03  & 2.13  & 97    & 214.26 & 0.99 \\
p07-t25-v2t2    & 11.58 & 107 & 1277.6  & 1.56  & 8.01  & 117   & 991.46 & 1.21 \\\hline
Average     & 15.85 & 82  & 689.62  & 0.98  & 13.18 & 81.13 & 582.59 & 0.83 \\\hline
\end{tabular}
\end{table}

To be able to see the effect of ESPPRC-CS in larger instances, we designed two more experiments
using instances with 40 customers. In the first design, first, we solved heuristic branch
and price algorithm (H-BP) described in Section \ref{sec:primal-feas}. In H-BP,
we applied 2p-ESPPRC-CS with subset size 15 for at most 20 iterations. Then we applied
2p-ESPPRC-LL with 5, 10 and 15 label limits. After running H-BP with a time limit of
2 CPU hours, we started exact branch and price (E-BP) using the upper bound and columns
obtained by H-BP. This two step algorithm is referred as  E-BP-2S  in Section (\ref{sec:alg-view}).
At each node of E-BP, we applied  2p-ESPPRC-CS with subset size 15 for at most 15 iterations.
After, algorithm 2p-ESPPRC-LL with 50, 100 and 200 label limits was applied and the evaluation
of the node terminated with exact pricing 2p-ESPPRC. The second step, E-BP was run with a time limit
of 6 CPU hours. To evaluate the contribution of ESPPRC-CS to the overall performance of the
algorithm, in the second experiment, we applied same 2 step branch and price algorithm
except in H-BP and E-BP we did not apply 2p-ESPPRC-CS. We refer the first branch and price
algorithms in the first design including algorithm 2p-ESPPRC-CS as H-BP-W and E-BP-W,
and the algorithms in the second design without 2p-ESPPRC-CS as H-BP-WO and E-BP-WO.
Table \ref{compexp-comp-pri3} shows the results of the two designs and presents comparisons.

In Table \ref{compexp-comp-pri3}, for each design (2 step branch and price algorithm with
and without 2p-ESPPRC-CS), we report upper bounds obtained at the end of H-BP-W and H-BP-WO,
and the optimality gaps at the end of E-BP-W and E-BP-WO. Therefore, we
can evaluate the effect of using 2p-ESPPRC-CS to both steps of the algorithm, H-BP and E-BP.
%in H-BP to the upper bound obtained, and effect of using 2p-ESPPRC-CS in E-BP to the gap at
%the end.
In addition, last three columns in Table \ref{compexp-comp-pri3} present the ratio of the
running times in the design with 2p-ESPPRC-CS to the times in the design without 2p-ESPPRC-CS
(denoted as {\em W/WO}). We report fraction of times in each of H-BP, E-BP and total
algorithm time (H-BP + E-BP).
%these two times in these last three columns.

Table \ref{compexp-comp-pri3}, indicates that addition of  2p-ESSPRC-CS to the algorithm
improves the performance of both H-BP and E-BP. In H-BP, in
14 out of 32 instances, the upper bounds obtained using H-BP-W (the algorithm with
2p-ESSPRC-CS) are strictly better while H-BP-WO
%the algorithm without 2p-ESSPRC-CS
provides better upper bounds in 3 out of 32 instances. In E-BP, the algorithm with
2p-ESSPRC-CS results strictly better optimality gaps in 13 instances while this number
is only 2 for the algorithm without 2p-ESSPRC-CS. In time comparison,
even though the time spent in H-BP for the algorithm with 2p-ESSPRC-CS is 1.6 times
greater than the corresponding time in the algorithm without 2p-ESSPRC-CS, in only 10 out of
32 instances time in H-BP for the algorithm with 2p-ESSPRC-CS is greater than the
the algorithm without 2p-ESSPRC-CS.
%The ratio 1.58 that is greater than 1 results because of the high ratio in instance p07-t40-v2t1.
2p-ESSPRC-CS results about \%30 improvement in both the time in E-BP and total time.
Therefore, we can conclude that the effect of heuristic pricing 2p-ESSPRC-CS for larger
instances is more significant and using 2p-ESSPRC-CS boosts the
overall performance of the algorithm, as well as the performance of H-BP and E-BP.

% are about \%30
%less in the  algorithm with 2p-ESSPRC-CS.
\begin{table}[!h]
\centering
\caption{Performance improvement with 2p-ESSPRC-CS in BP algorithm for 40 customers}\label{compexp-comp-pri3}
\begin {tabular}[h]{|c|c|c|c|c|c|c|c|}
\hline
\mbox{} & \multicolumn{2}{|c|}{E-BP-2S Without} & \multicolumn{2}{|c|}{E-BP-2S With} & \multicolumn{3}{|c|}{Time}\\
\mbox{} & \multicolumn{2}{|c|}{ESPPRC-CS} & \multicolumn{2}{|c|}{ESPPRC-CS} & \multicolumn{3}{|c|}{Comparison}\\\cline{2-8}
\mbox{}   & IP    & \% Gap   & IP      & \% Gap & H-BP  & E-BP & Tot Time \\
Data File &  H-BP & E-BP     & H-BP    & E-BP   & W/WO  & W/WO & W/WO \\\hline
p01-f40-v1t1 &  7131    & 0.01  & 7131  & 0 & 1.24  & 0.42  & 0.5   \\
p01-f40-v1t2 &  6986    & 0.14  & 7170  & 0.32  & 1 & 1 & 1 \\
p01-f40-v2t1 &  7017    & 0 & 6866  & 0 & 1.31  & 0.46  & 0.47  \\
p01-f40-v2t2 &  6821    & 1.99  & 6821  & 1.26  & 0.29  & 0.97  & 0.81  \\\hline
p01-l40-v1t1 &  7252    & 4.14  & 7183  & 3.22  & 1 & 0.99  & 0.99  \\
p01-l40-v1t2 &  6934    & 0 & 6934  & 0 & 0.64  & 0.49  & 0.5   \\
p01-l40-v2t1 &  6823    & 0 & 6823  & 0 & 0.33  & 0.68  & 0.67  \\
p01-l40-v2t2 &  6634    & 0 & 6633  & 0 & 1.26  & 0.58  & 0.6   \\\hline
p03-f40-v1t1 &  8780    & 0 & 8780  & 0 & 0.94  & 1.08  & 1.07  \\
p03-f40-v1t2 &  8756    & 0 & 8753  & 0 & 1.66  & 0.66  & 0.73  \\
p03-f40-v2t1 &  8663    & 0 & 8663  & 0 & 1.53  & 0.37  & 0.57  \\
p03-f40-v2t2 &  8456    & 0 & 8438  & 0 & 2.91  & 0.55  & 0.59  \\\hline
p03-m40-v1t1 &  7793    & 0.23  & 7740  & 0 & 1 & 1 & 1 \\
p03-m40-v1t2 &  7544    & 6.32  & 7236  & 2.95  & 0.7   & 0.97  & 0.9   \\
p03-m40-v2t1 &  7140    & 2.29  & 7132  & 2.79  & 1 & 0.97  & 0.98  \\
p03-m40-v2t2 &  7430    & 8.52  & 6947  & 4.18  & 0.31  & 0.97  & 0.81  \\\hline
p03-l40-v1t1 &  7167    & 0 & 7167  & 0 & 0.06  & 0.44  & 0.37  \\
p03-l40-v1t2 &  6950    & 0.07  & 6950  & 0.05  & 1 & 1 & 1 \\
p03-l40-v2t1 &  6857    & 0 & 6849  & 0 & 0.03  & 0.54  & 0.35  \\
p03-l40-v2t2 &  6846    & 3.19  & 6846  & 0 & 1.09  & 0.81  & 0.83  \\\hline
p07-f40-v1t1 &  7167    & 0 & 7167  & 0 & 0.37  & 0.5   & 0.48  \\
p07-f40-v1t2 &  7001    & 0 & 7001  & 0 & 0.29  & 0.55  & 0.49  \\
p07-f40-v2t1 &  6925    & 0.01  & 6881  & 0 & 0.33  & 0.41  & 0.41  \\
p07-f40-v2t2 &  6851    & 2.55  & 6845  & 2.34  & 1.41  & 1 & 1.01  \\\hline
p07-s40-v1t1 &  7343    & 0 & 7343  & 0 & 0.57  & 0.55  & 0.55  \\
p07-s40-v1t2 &  7117    & 0 & 7124  & 0 & 0.16  & 0.61  & 0.38  \\
p07-s40-v2t1 &  7000    & 0 & 7005  & 0 & 0.1   & 0.56  & 0.49  \\
p07-s40-v2t2 &  6944    & 0 & 6944  & 0 & 0.43  & 0.64  & 0.63  \\\hline
p07-t40-v1t1 &  7012    & 0.01  & 7013  & 0 & 0.6   & 0.39  & 0.4   \\
p07-t40-v1t2 &  6972    & 1.22  & 6972  & 0.35  & 2.43  & 0.99  & 1.1   \\
p07-t40-v2t1 &  6869    & 3.01  & 6857  & 0.04  & 24.3  & 0.81  & 0.94  \\
p07-t40-v2t2 &  6667    & 1.78  & 6658  & 1.64  & 0.22  & 1.01  & 0.81  \\\hline
Average      & \mbox{}  & \mbox{} & \mbox{} & \mbox{} & 1.58 &  0.72 &  0.7 \\\hline
\end{tabular}
\end{table}

\begin{comment}
\begin{table}[h]
\centering
 \begin {tabular}[t]{|ccccc|ccccc|}
Data & \multicolumn{2}{c}{{\bf With out SS}} &  \multicolumn{2}{c}{{\bf With SS}} & Data & \multicolumn{2}{c}{{\bf With out SS}} &  \multicolumn{2}{c}{{\bf With SS}} \\%\cline{2-10}
\mbox{} & {\bf UB-S1} & {\bf Gap-S2} &  {\bf UB-S1} & {\bf Gap-S2} &\mbox{} & {\bf UB-S1} & {\bf Gap-S2}  & {\bf UB-S1} & {\bf Gap-S2} \\\hline
a40-v1t1    &   7131    &   0.01\%  &   7131    &   0   &   e40-v1t1    &   7167    &   0   &   7167    &   0   \\
a40-v1t2    &{\bf 6986} &   0.14\%  &   7170    &   0.32\%  &   e40-v1t2    &   6950    &   0.07\%  &   6950    &  {\bf 0.05\%} \\
a40-v2t1    &   7017    &   0   &    {\bf 6866} &   0   &   e40-v2t1    &   6857    &   0   &   {\bf 6849}  &   0   \\
a40-v2t2    &   6821    &   1.99\%  &   6821    & {\bf 1.26\%}  &   e40-v2t2    &   6846    &   3.19\%  &   6846    &   {\bf 0} \\
b40-v1t1    &   7252    &   4.14\%  &    {\bf 7183} & {\bf 3.22\%}  &   f40-v1t1    &   7167    &   0   &   7167    &   0   \\
b40-v1t2    &   6934    &   0   &   6934    &   0   &   f40-v1t2    &   7001    &   0   &   7001    &   0   \\
b40-v2t1    &   6823    &   0   &   6823    &   0   &   f40-v2t1    &   6925    &   0.01\%  &   {\bf 6881}  &   {\bf 0} \\
b40-v2t2    &   6634    &   0   &    {\bf 6633} &   0   &   f40-v2t2    &   6851    &   2.55\%  &   {\bf 6845}  &   {\bf 2.34\%}    \\
c40-v1t1    &   8780    &   0   &   8780    &   0   &   g40-v1t1    &   7343    &   0   &   7343    &   0   \\
c40-v1t2    &   8756    &   0   &    {\bf 8753} &   0   &   g40-v1t2    &   {\bf 7117}  &   0   &   7124    &   0   \\
c40-v2t1    &   8663    &   0   &   8663    &   0   &   g40-v2t1    &   {\bf 7000}  &   0   &   7005    &   0   \\
c40-v2t2    &   8456    &   0   &    {\bf 8438} &   0   &   g40-v2t2    &   6944    &   0   &   6944    &   0   \\
d40-v1t1    &   7793    &   0.23\%  &    {\bf 7740} & {\bf 0}   &   h40-v1t1    &  {\bf 7012}   &   0.01\%  &   7013    &   {\bf 0} \\
d40-v1t2    &   7544    &   6.32\%  &   {\bf 7236}  & {\bf 2.95\%}  &   h40-v1t2    &   6972    &   1.22\%  &   6972    &   {\bf 0.35\%}    \\
d40-v2t1    &   7140    &   2.29\%  &   {\bf 7132}  & {\bf 2.79\%}  &   h40-v2t1    &   6869    &   3.01\%  & {\bf 6857}    &   {\bf 0.04\%}    \\
d40-v2t2    &   7430    &   8.52\%  &   {\bf 6947}  & {\bf 4.18\%}  &   h40-v2t2    &   6667    &   1.78\%  & {\bf 6658}    &       {\bf 1.64\%}
\\\hline
\end{tabular}
%\end{footnotesize}
%\end{scriptsize}
%\end{tiny}
\end{table}
\end{comment}

\subsection{Performance of Primal Feasible Heuristics}
\label{comp:prifeasHeur}
As described in Chapter 2, in our study, we could employ three approaches
to find primal feasible solutions for LRSP: construction base simple heuristic
I-Heur (described in Section \ref{sec:initcols}), a standard branch and bound using
MINTO and heuristic branch and price algorithm H-BP (described in Section
\ref{sec:primal-feas}). In this Section, we empirically demonstrate the quality of
the upper bound obtained using each heuristic approach.
Branch and bound approach may be used at the root node
of either heuristic branch and price (H-BP) or exact branch and price algorithm (E-BP).
Here, in this experiment, we used branch and bound (refer as BB) at the root node
of H-BP based on the RMP with the set of columns present at the end of root node.
We set a time limit of 1 CPU hour for BB. In the experiment, H-BP is run with a time
limit of 2 CPU hours and only heuristic pricing algorithm 2p-ESPPRC-LL with label
limits of 5 and 10  is used in the column generation algorithm.

Table \ref{primal-heur-table1} reports upper bound values obtained from and time spent
at each heuristic bounding procedure, I-Heur, BB and H-BP. Time for BB is the time to
solve the IP model and does not include the time required to solve the root node of H-BP
and generate the columns in the model to be used as input for BB algorithm.
I-Heur is very fast. Time required for BB and H-BP depends on the size of the instance;
they terminate very quickly in 25-customer instances, 40-customer instances require more time
and may terminate with the time limit. Even though, BB and H-BP are not very quick for
large instances, they improve the performance of the exact  branch and price algorithm
significantly by providing very good upper bounds for large instances.

Table  \ref{primal-heur-table2} reports the \% gaps between the known LP and optimal (or best)
solutions of each instance and upper bounds obtained from each heuristic bounding approach.
As seen from Table \ref{primal-heur-table2} I-Heur and BB provide upper bounds which are
about \%5 to \%11 away from the LP relaxation bounds. BB bounds are slightly better than the I-Heur
bounds. This difference is larger for 40 customer instances; BB provides bounds that are \%4.8
away from the LP relaxation bound on average while the average gap is 6.7 in I-Heur.
H-BP provides the strongest bound compared to I-Heur and BB.
The average gap between H-BP upper bounds and the LP relaxation bounds is
about 3.4 in all instances, and the average values are similar in both 25- and 40-customer
instances. Consistent results are obtained for the gaps from the known optimal or best
IP solutions. I-Heur provides bounds that are about \%4.3 away from the optimal
solution on average. The effect of BB is more significant for 40-customer instances
compared to 25-customer instances. H-BP is very strong in almost all instances;
in 15 out of 16 instances, H-BP bound is either optimal or at most \%1.5 away
from the optimal.

These results suggest that all of the bounding procedures perform well and provide
good upper bounds. Depending on the performance of the overall exact branch and price
algorithm, one or more of these heuristics can be employed. If the exact branch and price
algorithm terminates quickly, then using only I-Heur or BB may be enough.
\begin{table}[!h]
\centering
\caption{Performance of I-Heur, BB, H-BP}\label{primal-heur-table1}
\begin {tabular}[!h]{|c|c|c|c|c|c|c|}
\hline
\mbox{} & \multicolumn{3}{|c|}{Upper Bound Values} & \multicolumn{3}{|c|}{Time (s)} \\\cline{2-7}
Data file & I-Heur & BB in H-BP Root & H-BP & I-Heur & BB in H-BP Root & H-BP \\\hline
p01-f25-v2t2 &  4722 &  4710    &   4425    & 0.04  &   21.15   &   264.96  \\
p01-l25-v2t2 &  4762 &  4762    &   4650    & 0.04  &   0.48    &   25.32   \\
p03-f25-v2t2 &  4831 &  4831    &   4742    & 0.04  &   11.38   &   135.94  \\
p03-m25-v2t2 &  4767 &  4761    &   4548    & 0.03  &   6.87    &   80.35   \\
p03-l25-v2t2 &  4805 &  4794    &   4734    & 0.03  &   1.56    &   45.28   \\
p07-f25-v2t2 &  4721 &  4576    &   4483    & 0.03  &   0.09    &   71.32   \\
p07-s25-v2t2 &  4934 &  4856    &   4781    & 0.04  &   3.03    &   57.88   \\
p07-t25-v2t2 &  4769 &  4769    &   4654    & 0.03  &   1.71    &   40.63   \\\hline
Average & \mbox{}&  \mbox{} &   \mbox{} & 0.04  &   5.78    &   90.21   \\\hline
p01-f40-v1t1 &  7170 &  7170    &   7131    & 0.05  &   14.35   &   1452.48 \\
p01-l40-v1t1 &  7449 &  7252    &   7252    & 0.07  &   46.31   &   7201.3  \\
p03-f40-v1t1 &  9253 &  9009    &   8780    & 0.04  &   38.73   &   608.19  \\
p03-m40-v1t1 &  7827 &  7793    &   7793    & 0.06  &   2589.2  &   7201.11 \\
p03-l40-v1t1 &  7386 &  7386    &   7167    & 0.05  &   190.03  &   2222.19 \\
p07-f40-v1t1 &  7424 &  7324    &   7167    & 0.04  &   63.21   &   680.46  \\
p07-s40-v1t1 &  7860 &  7343    &   7343    & 0.05  &   0.28    &   75.01   \\
p07-t40-v1t1 &  7283 &  7283    &   7012    & 0.05  &   51.52   &   825.27  \\\hline
Average & \mbox{}    &  \mbox{} &   \mbox{} & 0.05  &   374.2   &   2533.25 \\\hline
\end{tabular}
\end{table}
\begin{table}[!h]
\centering
\caption{I-Heur, BB, H-BP Upper Bounds vs. LP Relaxations and Optimal IPs }\label{primal-heur-table2}
\begin {tabular}[!h]{|c|c|c|c|c|c|c|}
\hline
\mbox{} & \multicolumn{3}{|c|}{\%Gap from LP Relax.} & \multicolumn{3}{|c|}{\%Gap from Optimal (or Best IP)} \\\cline{2-7}
Data file & I-Heur & BB in H-BP Root & H-BP & I-Heur & BB in H-BP Root & H-BP \\\hline
p01-f25-v2t2    & 9.93  & 9.65  & 3.02  & 6.71  & 6.44  &   0   \\
p01-l25-v2t2    & 4.96  & 4.96  & 2.49  & 3.48  & 3.48  &   1.04    \\
p03-f25-v2t2    & 7.29  & 7.29  & 5.32  & 1.88  & 1.88  &   0   \\
p03-m25-v2t2    & 8.44  & 8.3   & 3.45  & 6.43  & 6.3   &   1.54    \\
p03-l25-v2t2    & 6.6   & 6.35  & 5.02  & 3.22  & 2.99  &   1.7 \\
p07-f25-v2t2    & 7.58  & 4.28  & 2.16  & 5.33  & 2.1   &   0.02    \\
p07-s25-v2t2    & 7.03  & 5.34  & 3.71  & 3.2   & 1.57  &   0   \\
p07-t25-v2t2    & 5.67  & 5.67  & 3.12  & 3.85  & 3.85  &   1.35    \\\hline
Average &   7.19    & 6.48  & 3.54  & 4.26  & 3.57  &   0.71    \\\hline
p01-f40-v1t1    & 2.78  & 2.78  & 2.22  & 0.55  & 0.55  &   0   \\
p01-l40-v1t1    & 7.96  & 5.11  & 5.11  & 3.7   & 0.96  &   0.96    \\
p03-f40-v1t1    & 6.39  & 3.58  & 0.95  & 5.39  & 2.61  &   0   \\
p03-m40-v1t1    & 11.01 & 10.53 & 10.53 & 7.91  & 7.45  &   7.45    \\
p03-l40-v1t1    & 6.73  & 6.73  & 3.56  & 3.06  & 3.06  &   0   \\
p07-f40-v1t1    & 5.26  & 3.84  & 1.62  & 3.59  & 2.19  &   0   \\
p07-s40-v1t1    & 8.32  & 1.19  & 1.19  & 7.04  & 0 &   0   \\
p07-t40-v1t1    & 4.93  & 4.93  & 1.03  & 3.86  & 3.86  &   0   \\\hline
Average &   6.67    & 4.84  & 3.28  & 4.39  & 2.58  &   1.05    \\\hline
\end{tabular}
\end{table}

Earlier in this section, we noted that BB approach may be used at the root node
of either H-BP or E-BP. In our experiments in Tables \ref{primal-heur-table1}
and \ref{primal-heur-table2}, we used BB at the root node of H-BP.
We also tested the performance of the E-BP algorithm in which BB is used
at the root node to provide an upper bound. Table \ref{primal-heur-table3}
presents the results comparing an E-BP with BB upper bounding procedure at root
node with an E-BP without BB procedure. Table \ref{primal-heur-table3}
reports the gap between the upper bound obtained using BB and the
known optimal IP solution, number of evaluated branch and price nodes
in both E-BP with BB and E-BP without BB, and the fraction of total
algorithm time of E-BP with BB to the time of E-BP without BB.
For both 25- and 40-customer instances, total algorithm time increases
if BB is employed at root node. On average, algorithm time increases
about \%13.  The number of evaluated nodes is decreased
in only 5 out of 16 instances with the BB procedure. Even though
the upper bound obtained using BB procedure is about \%2.5 away
from the optimal solution, using BB at the root node of E-BP does not
improve, but worsen the overall algorithm performance.
Therefore, based on these results, we have decided to use BB only at the
root node of H-BP.

\begin{table}[!h]
\centering
\caption{Performance of BB at Root Node of E-BP}\label{primal-heur-table3}
\begin {tabular}[!h]{|c|c|c|c|c|}
\hline
\mbox{}   & \multicolumn{2}{|c|}{E-BP, BB at Root} & E-BP & Time \\\cline{2-3}
Data File & BBIP-Opt Gap & Nds & Nds & BB/WOBB \\\hline
p01-f25-v2t2 &  1.04    & 83    & 81    & 1.12  \\
p01-l25-v2t2 &  3.48    & 37    & 37    & 1.23  \\
p03-f25-v2t2 &  1.88    & 50    & 50    & 1.34  \\
p03-m25-v2t2 &  5.4 & 9 & 9 & 1.25  \\
p03-l25-v2t2 &  3.22    & 43    & 43    & 1.07  \\
p07-f25-v2t2 &  0.98    & 205   & 251   & 0.98  \\
p07-s25-v2t2 &  0   & 105   & 92    & 1.03  \\
p07-t25-v2t2 &  2.03    & 157   & 149   & 1.02  \\\hline
Average      &  2.25    & \mbox{} & \mbox{} & 1.13 \\\hline
p01-f40-v1t1 &  0.52    & 431   & 415   & 1.2   \\
p01-l40-v1t1 &  0.1 & 893   & 911   & 1 \\
p03-f40-v1t1 &  2.72    & 65    & 79    & 1 \\
p03-m40-v1t1 &  8.01    & 1199  & 1340  & 1 \\
p03-l40-v1t1 &  0.6 & 193   & 196   & 1.05  \\
p07-f40-v1t1 &  2.54    & 95    & 95    & 1.22  \\
p07-s40-v1t1 &  3.19    & 73    & 73    & 1.49  \\
p07-t40-v1t1 &  3.58    & 229   & 229   & 1.1   \\\hline
Average      &  2.66    & \mbox{}& \mbox{} & 1.13 \\\hline
\end{tabular}
\end{table}

\subsection{Performance of Initializing Branch \& Price Algorithm Using H-BP (E-BP vs E-BP-2S)}
As explained in  Section \ref{sec:BP-overview}, the upper bounding
procedures, I-Heur and H-BP, have also a role of providing a starting
set of feasible columns to initialize the branch and price algorithm.
When we solve exact branch and price algorithm, E-BP, we initialize it with
the columns and upper bound from I-Heur. Since H-BP is itself also a branch and
price algorithm, when we use H-BP to initialize exact branch and price algorithm,
E-BP, we call the overall algorithm as 2-step branch and price algorithm E-BP-2S.
The first step refers to H-BP, and the second step is the exact algorithm, E-BP.
As seen from the results in Section \ref{comp:prifeasHeur}, H-BP may take
too much time and this may worsen the total algorithm time. Even though H-BP may be
time consuming, it provides very strong upper bounds and spending this amount of time
for upper bound and nice columns will be worthwhile for large instances especially for
ones in which exact branch and price algorithm terminates with the time limit.
Here, we provide more computational results demonstrating the effect of
employing H-BP to initialize E-BP  to the overall branch and price algorithm.
In other words, we compare 1 step and 2 step algorithms, E-BP and E-BP-2S.

We designed two experiments. In the first experiment, we used 1 step branch and price
algorithm, i.e., we used I-Heur to initialize the E-BP. In E-BP, at each node
of the tree, we first employed heuristic pricing algorithm 2p-ESPPRC-LL with label limits
50, 100 and 200, and finally exact pricing algorithm 2p-ESPPRC.
We set total time limit to 8 CPU hours. In the second experiment, we used
2 step branch and price algorithm, i.e., we solved H-BP first. In H-BP, we
employed heuristic pricing algorithm 2p-ESPPRC-LL with label limits
of 5 and 10. We set time limit to 2 CPU hours. Then, E-BP started with the
columns and the upper bound from H-BP. Time limit was set to 6 CPU hours.

Table \ref{hbp-table25} and \ref{hbp-table40} present results of the
designed experiments for 25-customer and 40-customer instances, respectively.
Table \ref{hbp-table25} and \ref{hbp-table40} reports the number of evaluated
nodes in both E-BP and E-BP-2S and the fraction of total algorithm time
in E-BP to the E-BP-2S. In 25-customer instances, both algorithms
terminated with optimal solutions. However, in some of 40-customer instances
the algorithms terminate with optimality gap. Thus, Table \ref{hbp-table40}
also reports \% gap at the end of the algorithm. Table \ref{hbp-table25}
indicates that on average, 2 step algorithm is about \%18 faster than 1 step E-BP.
Initializing the exact branch and price algorithm with the H-BP
boosts the overall performance of the algorithm. There are also some instances
for which E-BP performs better than E-BP-2S. For most of the 25-customer
instances, computation times are very short. The average computation time
for instances in  Table \ref{hbp-table25} is about 340 CPU seconds
in E-BP and about 295 seconds in E-BP-2S. Therefore, using 1 step exact branch
and price algorithm, E-BP for 25-customer instances is not a bad
choice. We may prefer using 1 step algorithm for this size of instances.

Table \ref{hbp-table40} reports that on average E-BP time is \%8 better than
E-BP-2S time. However, Table \ref{hbp-table40} shows that only 6 out of 16 instances
E-BP time is strictly better than E-BP-2S while E-BP-2S is strictly better
in 7 instances. In addition, when we compare the optimality gap columns for
two designs, we can see that E-BP-2S is more successful in closing the gaps.
The strong upper bound obtained in H-BP has a significant role in the termination
with better gaps. Therefore, for large instances, since the algorithm
terminates with the time limit, for the sake of good optimality gaps, we prefer
using E-BP-2S.

\begin{table}[!h]
\centering
\caption{E-BP vs E-BP-2S: Effect of H-BP, 25-Customer Instances}\label{hbp-table25}
\begin {tabular}[!h]{|c|c|c|c|c|c|}
\hline
\mbox{}   & \multicolumn{2}{|c|}{Nds} & \multicolumn{2}{|c|}{CPU Time (s)} & Time  \\\cline{2-5}
Data File & E-BP  & E-BP-2S & E-BP  & E-BP-2S & E-BP/E-BP-2S \\\hline
p01-f25-v1t1 &  25   &  25   &  77.81   &  64.26   &    1.21    \\
p01-f25-v2t2 &  119  &  115  &  954.41  &  665.3   &    1.43    \\
p01-l25-v1t1 &  24   &  23   &  35.45   &  29.14   &    1.22    \\
p01-l25-v2t2 &  39   &  35   &  414.51  &  345.96  &    1.2 \\
p03-f25-v1t1 &  236  &  231  &  182.13  &  227.96  &    0.8 \\
p03-f25-v2t2 &  52   &  43   &  316.88  &  233.05  &    1.36    \\
p03-m25-v1t1 &  509  &  483  &  314.12  &  431.17  &    0.73    \\
p03-m25-v2t2 &  9    &  9    &  523.75  &  284.47  &    1.84    \\
p03-l25-v1t1 &  17   &  15   &  36.21   &  29.6    &    1.22    \\
p03-l25-v2t2 &  37   &  45   &  452.17  &  453.74  &    1   \\
p07-f25-v1t1 &  55   &  37   &  81.28   &  63.86   &    1.27    \\
p07-f25-v2t2 &  197  &  135  &  907.33  &  619.14  &    1.47    \\
p07-s25-v1t1 &  599  &  197  &  87.77   &  52.2    &    1.68    \\
p07-s25-v2t2 &  100  &  93   &  217.97  &  179.56  &    1.21    \\
p07-t25-v1t1 &  35   &  47   &  51.8    &  44.94   &    1.15    \\
p07-t25-v2t2 &  91   &  211  &  754.9   &  997.57  &    0.76    \\\hline
average & \mbox{}   & \mbox{} & 338.03  &  295.12  &    1.22    \\\hline
\end{tabular}
\end{table}

\begin{table}[!h]
\centering
\caption{E-BP vs E-BP-2S: Effect of H-BP}\label{hbp-table40}
\begin {tabular}[!h]{|c|c|c|c|c|c|}
\hline
\mbox{}   & \multicolumn{2}{|c|}{Gap \%} & \multicolumn{2}{|c|}{Nds} & Time  \\\cline{2-5}
Data File & E-BP  & E-BP-2S & E-BP  & E-BP-2S &  E-BP/E-BP-2S \\\hline
p01-f40-v1t1 &  0    &  0.01    & 415   & 909   & 0.19  \\
p01-f40-v2t2 &  4.28 &  1.99    & 15    & 16    & 1 \\
p01-l40-v1t1 &  6.9  &  4.14    & 911   & 324   & 0.99  \\
p01-l40-v2t2 &  0    &  0   & 25    & 25    & 1.01  \\
p03-f40-v1t1 &  0    &  0   & 79    & 65    & 0.52  \\
p03-f40-v2t2 &  0    &  0   & 11    & 11    & 1.19  \\
p03-m40-v1t1 &  0.1  &  0.23    & 1340  & 1021  & 1 \\
p03-m40-v2t2 &  8.07 &  8.52    & 32    & 24    & 1 \\
p03-l40-v1t1 &  0    &  0   & 196   & 197   & 0.86  \\
p03-l40-v2t2 &  5.53 &  3.19    & 20    & 22    & 1.22  \\
p07-f40-v1t1 &  0    &  0   & 95    & 101   & 0.86  \\
p07-f40-v2t2 &  4.55 &  2.55    & 34    & 19    & 1.28  \\
p07-s40-v1t1 &  0    &  0   & 73    & 65    & 1.14  \\
p07-s40-v2t2 &  0    &  0   & 213   & 173   & 1.12  \\
p07-t40-v1t1 &  0    &  0.01    & 229   & 220   & 0.25  \\
p07-t40-v2t2 &  5.71 &  1.78    & 6 & 5 & 1.04  \\\hline
Average & \mbox{}    &  \mbox{} & \mbox{}& \mbox{}& 0.92 \\\hline
\end{tabular}
\end{table}


\subsection{Results of Branch \& Price Algorithm for Generated Instances}

In this section, we present solution and performance details of the exact branch
and price algorithms for 25- and 40-customer instances. Based on the results of our
previous experiments, we came up with ideal designs of branch and price
algorithms to solve the generated LRSP instances. The objective of the algorithms
is to provide optimal or best solutions, and optimality gaps in case of
not proving optimality for the LRSP instances. Here, we present computational results
for all 64 instances in 25- and 40-customer instance sets.
%We give optimal
%(or best) solution's details for all instances.
%and described in Section \ref{sec:BP-overview}.

\paragraph{25-Customer Instances}
Our previous computational experiments showed that our 1 step exact branch and price
algorithm, E-BP is as successful as our 2 step exact branch and price algorithm,
E-BP-2S for medium-size instances. Therefore, for 25-customer instances,
we may either use E-BP or E-BP-2S. Here,  we used E-BP which only
employs I-Heur to initialize.
%In this section, we only use I-Heur to initialize the exact branch and price algorithm.
%And in the next section, we will provide results for 2-step algorithm.
%Thus, this and the next section provide a preformance comparison of whether
%using initialization H-BP or not.

Based on  the performance results from our previous experiments, we designed the
following exact branch and price algorithm, E-BP for 25-customer instances.
The main components and input parameters of E-BP are as follows:
\begin{itemize}
\item Run I-Heur to find initial columns and an initial upper bound.
\item At each node node of the tree, column generation algorithm
  uses the following pricing algorithms:
  \begin{itemize}
  %\item Run at most 12 iterations 2p-ESPPRC-CS(12) with subset size 12.
  \item Run 2p-ESPPRC-LL with label limit 50, 100 and 200.
    The main
    implementation steps are described in Section \ref{pri-E-BP}.
\zaupdate{may move the implementation details from ch2 to here.}
  \item Run 2p-ESPPRC.
  \end{itemize}
\item Set time limit of 6 CPU hours.
\end{itemize}

Table \ref{25Cust_res} shows the results for instances  with 25 customers.
For each instance, Table \ref{25Cust_res} reports the upper bound obtained from I-Heur (I UB),
the LP relaxation value at the root node (LP), the optimal objective value (IP),
the integrality gap at the end of the running time (Gap \%), the number of evaluated
nodes (Nds), and the total CPU time in minutes (CPU). The last three columns present
some details about the optimal solution: the indices of open facilities
(Loc), the number of vehicles used at each open facility (\# Veh), and the total number
of routes in each solution (\# Rt).

Overall, the algorithm handled the 25-customer instances well; 29 of the
32 instances were solved to optimality in less than 25 minutes with the majority (24)
solved in 10 minutes or less. For only one instance did the algorithm terminate due to the
time limit before proving optimality. The initial feasible solutions provided by I-Heur
were good for the 25-customer instances with an average deviation of 4.4\% from optimality
which is consistent with the results from Section \ref{comp:prifeasHeur}.
The LP relaxation bounds were quite strong, as expected with the addition of valid inequalities
at the root node. The average deviation of the LP relaxation bound from the optimality
was 2.6\%, which is in line with the results of the smaller instances (instances with 20 customers
in Section \ref{alternativeForm}).

In addition to the solution summary given in the last three columns of Table \ref{25Cust_res},
we listed the constructed vehicle routes and schedules (i.e., selected pairings) in the optimal
solutions for all instances in Appendix \ref{app:25customerSolutions}. In all instances, the
number of facilities selected to be used is 2. In addition, in Table \ref{SolDet}, we present
some summary statistics about the solutions of all 32 instances, such as the number of vehicles
used, the average number of routes per vehicle, the average number of customers per route
and per vehicle, and the maximum number of customers per vehicle. In all instances, total number
of vehicles used changes between 3 to 5 while number of vehicles per facility changes between
1 to 3. Again, based on the table, we can see that each vehicle's schedule covers at least 2
routes (2.1 routes per vehicle on average) and some vehicles cover more than 3 routes
(average \# of routes is between 1.5 to 3.3), which may include 11 customers at most.
These statistics give us an intuition about the difficulty of the instances.

\begin{table}[!htbp]
\centering
\caption{Results for instances with 25 customers and 5 facilities}\label{25Cust_res}
\begin {tabular}[!h]{|c|c|c|c|c|r|r|c|c|c|}
\hline
Data File & I UB & LP & IP$^1$ & Gap \% & Nds & CPU (m) & Loc & \# Veh & \# Rt \\ \hline
p01-f25-v1t1 &  4795 &  4504.4 & 4699   & 0 & 25    & 1.3   & 0,4   & 2,2   & 10 \\
p01-f25-v1t2 &  4795 &  4425.8 & 4511   & 0 & 297   & 208.1 & 0,4   & 1,2   & 10 \\
p01-f25-v2t1 &  4722 &  4366.1 & 4425   & 0 & 99    & 9.9   & 0,4   & 1,2   & 7 \\
p01-f25-v2t2 &  4722 &  4295.4 & 4425   & 0 & 119   & 15.9  & 0,4   & 1,2   & 7 \\\hline
p01-l25-v1t1 &  5193 &  4849.8 & 4899   & 0 & 24    & 0.6   & 0,1   & 2,2   & 8 \\
p01-l25-v1t2 &  4968 &  4725.4 & 4815   & 0 & 139   & 2.9   & 0,4   & 2,2   & 8 \\
p01-l25-v2t1 &  4809 &  4669.0 & 4759   & 0 & 157   & 5.7   & 0,4   & 2,2   & 6 \\
p01-l25-v2t2 &  4762 &  4537.1 & 4602   & 0 & 39    & 6.9   & 0,1   & 1,2   & 7 \\\hline
p03-f25-v1t1 &  5281 &  4810.0 & 5062   & 0 & 236   & 3.0   & 0,2   & 3,2   & 10 \\
p03-f25-v1t2 &  5139 &  4701.1 & 4837   & 0 & 77    & 2.5   & 0,2   & 2,2   & 10 \\
p03-f25-v2t1 &  5048 &  4606.1 & 4758   & 0 & 393   & 4.6   & 0,4   & 2,2   & 8 \\
p03-f25-v2t2 &  4831 &  4502.5 & 4742   & 0 & 52    & 5.3   & 0,4   & 2,2   & 8 \\\hline
p03-m25-v1t1 &  4995 &  4697.4 & 4785   & 0 & 509   & 5.2   & 0,4   & 2,2   & 10 \\
p03-m25-v1t2 &  4869 &  4578.2 & 4785   & 0.1 & 3790    & 360.0 & 0,4   & 2,2   & 10 \\
p03-m25-v2t1 &  4992 &  4485.8 & 4689   & 0 & 59    & 2.9   & 0,4   & 2,2   & 7 \\
p03-m25-v2t2 &  4767 &  4396.1 & 4479   & 0 & 9         & 8.7   & 0,4   & 1,2   & 7 \\\hline
p03-l25-v1t1 &  5186 &  4795.3 & 4842   & 0 & 17    & 0.6   & 0,2   & 2,2   & 9 \\
p03-l25-v1t2 &  4871 &  4670.2 & 4803   & 0 & 37    & 1.9   & 0,4   & 2,2   & 7 \\
p03-l25-v2t1 &  4805 &  4620.9 & 4741   & 0 & 41    & 2.5   & 0,4   & 2,2   & 6 \\
p03-l25-v2t2 &  4805 &  4507.6 & 4655   & 0 & 37    & 7.5   & 2,4   & 1,2   & 7 \\\hline
p07-f25-v1t1 &  4870 &  4642.6 & 4761   & 0 & 55    & 1.4   & 0,2   & 1,3   & 7 \\
p07-f25-v1t2 &  4823 &  4544.5 & 4685   & 0 & 67    & 131.2 & 2,4   & 2,1   & 9 \\
p07-f25-v2t1 &  4721 &  4485.6 & 4703   & 0 & 229   & 8.8   & 0,2   & 2,2   & 6 \\
p07-f25-v2t2 &  4721 &  4388.3 & 4482   & 0 & 197   & 15.1  & 0,2   & 1,2   & 6 \\\hline
p07-s25-v1t1 &  5412 &  5000.1 & 5098   & 0 & 599   & 1.5   & 0,4   & 2,3   & 9 \\
p07-s25-v1t2 &  5238 &  4852.2 & 4889   & 0 & 123   & 2.9   & 0,4   & 2,2   & 9 \\
p07-s25-v2t1 &  5159 &  4739.4 & 4787   & 0 & 21    & 0.6   & 0,4   & 2,2   & 7 \\
p07-s25-v2t2 &  4934 &  4609.8 & 4781   & 0 & 100   & 3.6   & 0,4   & 2,2   & 8 \\\hline
p07-t25-v1t1 &  5134 &  4836.5 & 4898   & 0 & 35    & 0.9   & 0,1   & 2,2   & 8 \\
p07-t25-v1t2 &  5162 &  4703.1 & 4886   & 0 & 1345  & 25.0  & 0,2   & 2,2   & 8 \\
p07-t25-v2t1 &  5025 &  4646.2 & 4767   & 0 & 777   & 18.8  & 0,2   & 2,2   & 6 \\
p07-t25-v2t2 &  4769 &  4513.1 & 4592   & 0 & 91    & 12.6  & 0,1   & 1,2   & 6 \\\hline
\multicolumn{10}{l}{$^1$Optimal or best IP found by the algorithm within the time limit.}\\
\end{tabular}
\end{table}

\begin{table}[h!tbp]
\centering
\caption{More Solution Details for 25-customer instances}\label{SolDet}
\begin {tabular}[!h]{|l|c|c|} \hline
%\multicolumn{1}{|c}{} & \multicolumn{2}{|c||}{25-customer instances} \\\cline{2-3}
Case                          & Changes btw [min - max] & Avg.  \\\hline
\# of vehicles per facility      &  1 - 3 &  \mbox{}  \\
Total \# of vehicles             &  3 - 5  &  \mbox{}  \\
Avg. \# of routes per vehicle    &  1.5 - 3.3 & 2.1  \\
Avg. \# of customers per route   &  2.5 - 4.2 & 3.3  \\
Avg. \# of customers per vehicle &  5 - 8.3   & 6.9  \\
Max \# of customers per vehicle  &  7 - 11    & 8.5  \\\hline
\end{tabular}
\end{table}

As can be understood from the labels of the instances, the computational tests
can also measure the sensitivity of the algorithm to vehicle capacity and time
limit. As the labels of the instances indicate that there are 16 pairs of instances
which differ only in vehicle capacity (instances with label x-v1t1 and x-v2t1)
and 16 pairs which differ only in distance limits (the instances with label x-v1t1
and x-v1t2). And there are 8 instances in which both vehicle capacity
and time limit increases from a reference instance (in an instances with label
x-v2t2, both the vehicle capacity and the distance limit are higher then
the instance x-v1t1).
%For example, the instances with label x-v1t1 and x-v2t1
%differ only in vehicle capacity;   the instances with label x-v1t1 and x-v1t2
%differ only in distance limit; and in the instances with label x-v2t2 the
%both the vehicle capacity and the distance limit increases with respect to
%the instance x-v1t1.
As expected, since the pricing problem gets more difficult, the solution time
tends to increase with the vehicle capacity and the vehicle time limit. In addition,
the optimal cost does not increase when the vehicle capacity and/or the vehicle
time limit increase(s). In 10 out of 16 pairs the computation time increases
when the vehicle capacity increases, in 14 out of 16 pairs the computation time
increases when the distance limit increase. In 8 cases out of 8, the computation time
increases when both time and vehicle limits increase.
However, there are also instances which favor the increase in time limit and/or capacity.

An increase in the vehicle capacity and time limit can reduce the cost through a change
in the routing cost together with a decrease in the number of vehicles and/or a change
in the set of selected facilities.
%This can be seen through the solution details presented
%in Table  \ref{25Cust_res} and in Appendix \ref{app:25customerSolutions}.
On average, we calculated about \%3.7 and \%2.2 reduction in cost
with vehicle capacity increase and time limit increase, respectively. The cost
reduces about \%5.8 on average if both vehicle and time limit increases.
%We summarize some of the results in Table \ref{EffectVCTL}. The third
%column ($L^V$, $L\uparrow$) presents the cases where time limit increases for fixed vehicle
%capacity, the fourth column ($L^V \uparrow$, $L$) presents the cases where vehicle
%capacity increases for fixed time limit, and  the fifth column ($L^V \uparrow$, $L\uparrow$)
%presents the cases both time limit and vehicle capacity increase.
%Table \ref{EffectVCTL} shows  the number of cases in which total CPU time increased,
%the average increase in time with respect to the reference, and the average percentage
%decrease in cost with the corresponding change in vehicle capacity or time limit.
%25-customer instances seem to be sensitive in terms of computation times to the changes in
%vehicle capacity and time limit.

\begin{comment}
The solution information provided in Table \ref{25Cust_res} and Table \ref{40Cust_res2} highlights
the value of solving the integrated problem instead of a series of problems.  Although the same number
of facilities is selected within an instance group (meaning a fixed set of customer locations and
a set of candidate facility locations), the particular facilities selected may differ depending
on the vehicle capacity and time limit values.  Which facilities are selected impacts which customers
are served by each facility and hence the sequence of the routes and the structure of the pairings.
From the solution information, we can see that, within an instance group, increases in the vehicle capacity
and the time limit reduce the total cost through changes in facility, decreases in the number of vehicles
required, changes in the underlying route sequences.
\end{comment}

\begin{comment}
\begin{table}[h!tbp]
\centering
\caption{More Solution Detail}\label{SolDet}
\begin {tabular}[t]{|l|c|c||c|c|} \hline
\multicolumn{1}{|c}{} & \multicolumn{2}{|c|}{25-customer instances} & \multicolumn{2}{c|}{40-customer instances} \\\cline{2-5}
                          & min - max & Avg. &  min - max & Avg. \\\hline
\# of vehicles per facility      &  1 - 3 &  & 1 - 3 &  \\
Total \# of vehicles             &  3 - 5  &  & 4 - 6 &  \\
Avg. \# of routes per vehicle    &  1.5 - 3.3 & 2.1 & 1.8 - 3.2 & 2.5 \\
Avg. \# of customers per route   &  2.5 - 4.2 & 3.3 & 2.5 - 4.4 & 3.4 \\
Avg. \# of customers per vehicle &  5 - 8.3   & 6.9 & 6.7 - 10 & 8.3 \\
Max \# of customers per vehicle  &  7 - 11    & 8.5 & 7 - 14 & 10.2 \\\hline
\end{tabular}
\end{table}


\begin{table}[b!]
\centering
\caption{Effect of Vehicle Capacity and Time Limit for 25 customer instances}\label{EffectVCTL}
\begin {tabular}[t]{l|c|c|c}
Case & $L^V$, $L^T\uparrow$ & $L^V \uparrow$, $L^T$ & $L^V \uparrow$, $L^T\uparrow$ \\ \hline
\# of times Total CPU increases                      & 14 / 16 & 10 / 16 & 8 / 8 \\
 Avg. increase in CPU      & 34.9 & 4.7 & 11.6 \\
Avg. \% decrease in cost  & 3.7 & 2.2 & 5.9 \\\hline
\end{tabular}
\end{table}
\end{comment}

\begin{comment}
\begin{table}[b!]
\centering
\caption{Effect of Vehicle Capacity and Time Limit for 25 customer instances}\label{EffectVCTL}
\begin {tabular}[t]{l|l|c|c|c}
\mbox{} & \mbox{} & $L^V$, $L\uparrow$ & $L^V \uparrow$, $L$ & $L^V \uparrow$, $L\uparrow$ \\ \hline
25-customer  & \# of times Total CPU increases                      & 14 / 16 & 10 / 16 & 8 / 8 \\
\mbox{}      & Avg. increase in CPU      & 34.9 & 4.7 & 11.6 \\
\mbox{}      & Avg. \% decrease in cost  & 3.7 & 2.2 & 5.9 \\\hline
40-customer  & \#  of times Total CPU increases                       & 12 / 16 & 7 / 16 & 6 / 8 \\
\mbox{}      & Avg. increase in CPU      & 2.7 & 4.8 & 8.7 \\
\mbox{}      & Avg. \% decrease in cost  & 1.6 & 3.8 & 5.3 \\\hline
\end{tabular}
\end{table}
\end{comment}

%\subsection{2 Step Branch  \& Price Algorithm for  40 Customer Instances}
%\subsection{Performance of 2 Step Branch \& Price Algorithm: E-BP-2S}
\paragraph{40-Customer Instances}
Since our previous results suggested that 2 step branch and price algorithm
performs better for larger instances, here we used 2 step branch and price
algorithm (E-BP-2S) for 40-customer instances. As explained before, E-BP-2S
calls H-BP to obtain an initial set of columns and an upper bound, then
starts running E-BP with the obtained initialization from H-BP.
%In this section, we evaluated the performance of the 2 step branch-and-price
%algorithm described in Section \ref{sec:BP-overview} on the 40-customer instances.
Based on  the performance results from our previous experiments, we designed
%2 step exact branch and price algorithm referred as
E-BP-2S with the following steps and design parameters.
\begin{itemize}
\item Run H-BP with a time limit of 2 CPU hours. Steps of the H-BP are:
  \begin{itemize}
  \item Run I-Heur to find initial columns and an initial upper bound.
  \item At each node node of the tree, column generation algorithm
    uses the following pricing algorithms:
    \begin{itemize}
    \item Run at most 20 iterations 2p-ESPPRC-CS(15) with subset size 15.
    \item Run 2p-ESPPRC-LL with label limit 5 and 10.
    \end{itemize}
  \item At root node, run standard branch and bound (BB) with current
    set of columns to update the primal solution.
  \end{itemize}
\item Run E-BP with a time limit of 6 CPU hours. Steps of the E-BP are:
  \begin{itemize}
  \item Initialize E-BP with the columns and upper bound from H-BP.
  \item At each node node of the tree, column generation algorithm
    uses the following pricing algorithms:
    \begin{itemize}
    \item  Run 2p-ESPPRC-LL with label limit 50, 100 and 200. The main
      implementation steps are described in Section \ref{pri-E-BP}.
    \item Run 2p-ESPPRC.
    \end{itemize}
  \end{itemize}
\end{itemize}
Table \ref{40Cust_res2} shows the results for instances  with 40 customers.
For each instance, Table \ref{40Cust_res2} reports the upper bound obtained from
I-Heur (I UB), the upper bound obtained at the end of H-BP (H-BP UB),
the LP relaxation value at the root node (LP), the optimal objective value (IP),
the integrality gap at the end of the running time (Gap \%), the number of evaluated
nodes (Nds), the CPU times spent in H-BP, E-BP and the total CPU time in minutes.
%For the 40-customer instances,
%we used E-BP-2S with a time limit of 2 CPU hours for
%stage one and 6 CPU hours for stage two and reported the results in Table \ref{40Cust_res2}.
%The labels in Table \ref{40Cust_res2} are same as the ones in Table \ref{25Cust_res}
%but with three additional columns: the fifth column (H-BP UB) contains
%the upper bound obtained at the end of H-BP, the tenth and eleventh columns contain
%the CPU times spent in H-BP and E-BP, respectively.
As expected,
the solution times in the 40-customer instances were larger than the 25-customer instances.
Even though the algorithm terminated with a nonzero gap in some cases,
the algorithm was able to prove optimality most of the time.
For other instances, good solutions with small integrality gaps were obtained
within the time limit.  The gap was at most 0.35\% for four of the instances and between
1.26\% and 4.18\% for the other seven instances.

%The initial feasible solutions provided by I-Heur were good for the 25-customer
%instances with an average deviation of 4.4\% from optimality. Similarly,
%for the 40-customer instances,
The average deviation of the upper bound obtained from I-Heur from the optimal or best integer
solution was 4.5\%. However, the overall success of the algorithm can be
attributed to the effectiveness of H-BP in finding good upper bounds. In all but one
instance, the H-BP upper bound was better than the I-Heur upper bound, and for 23 instances
out of 32, the H-BP upper bound value equaled that of the optimal (or best integer) solution.
The LP relaxation bounds were quite strong at the root node; the average deviation of the
LP relaxation bound from the optimality was 2.2\% in 40-customer instances, which are in line with
the results of the smaller instances (instances with 20 customers in Section \ref{alternativeForm}).

\begin{table}[!htbp]
\centering
\caption{Results for instances with 40 customers and 5 facilities}\label{40Cust_res2}
\begin {tabular}[!h]{|c|c|c|c|c|r|r|r|r|r|}
\hline
\mbox{}   & I  & H-BP &    &        &    Gap     &       & \multicolumn{3}{|c|}{CPU(m)}\\\cline{8-10}
Data file & UB & UB   & LP & IP$^1$ & \% \ \ & Nds & H-BP & E-BP & Total \\ \hline
p01-f40-v1t1 &  7170 &  7131 &  6976.0  & 7131  & 0 & 807   & 53.2  & 150.0 & 203.2 \\
p01-f40-v1t2 &  7170 &  7170 &  6826.2  & 6999  & 0.32  & 1065  & 120.0 & 360.0 & 480.0 \\
p01-f40-v2t1 &  7064 &  6866 &  6700.2  & 6866  & 0 & 203   & 4.9   & 113.0 & 117.9 \\
p01-f40-v2t2 &  6984 &  6821 &  6581.0  & 6821  & 1.26  & 35    & 35.4  & 360.7 & 396.1 \\\hline
p01-l40-v1t1 &  7449 &  7183 &  6899.7  & 7183  & 3.22  & 573   & 120.1 & 360.4 & 480.5 \\
p01-l40-v1t2 &  7236 &  6934 &  6757.4  & 6934  & 0 & 43    & 5.6   & 48.0  & 53.6  \\
p01-l40-v2t1 &  7166 &  6823 &  6616.7  & 6823  & 0 & 129   & 2.4   & 103.2 & 105.6 \\
p01-l40-v2t2 &  7166 &  6633 &  6513.6  & 6612  & 0 & 25    & 8.4   & 187.6 & 196.0 \\\hline
p03-f40-v1t1 &  9253 &  8780 &  8697.4  & 8780  & 0 & 97    & 9.6   & 97.8  & 107.4 \\
p03-f40-v1t2 &  9028 &  8753 &  8578.6  & 8753  & 0 & 41    & 28.0  & 154.6 & 182.6 \\
p03-f40-v2t1 &  8943 &  8663 &  8439.0  & 8440  & 0 & 3 & 7.8   & 8.8   & 16.6  \\
p03-f40-v2t2 &  8943 &  8438 &  8390.4  & 8438  & 0 & 9 & 11.5  & 138.2 & 149.7 \\\hline
p03-m40-v1t1 &  7827 &  7740 &  7050.9  & 7253  & 0 & 2524  & 120.0 & 360.1 & 480.1 \\
p03-m40-v1t2 &  7544 &  7236 &  6907.7  & 7236  & 2.95  & 24    & 84.0  & 360.4 & 444.4 \\
p03-m40-v2t1 &  7479 &  7132 &  6793.9  & 7132  & 2.79  & 27    & 120.0 & 360.1 & 480.1 \\
p03-m40-v2t2 &  7430 &  6947 &  6655.5  & 6947  & 4.18  & 4 & 37.2  & 360.3 & 397.5 \\\hline
p03-l40-v1t1 &  7386 &  7167 &  6920.4  & 7167  & 0 & 183   & 1.5   & 55.6  & 57.1  \\
p03-l40-v1t2 &  7212 &  6950 &  6771.6  & 6950  & 0.05  & 2121  & 120.0 & 360.0 & 480.0 \\
p03-l40-v2t1 &  7121 &  6849 &  6662.9  & 6849  & 0 & 159   & 3.9   & 113.8 & 117.7 \\
p03-l40-v2t2 &  7018 &  6846 &  6544.7  & 6639  & 0 & 53    & 39.0  & 290.9 & 329.9 \\\hline
p07-f40-v1t1 &  7424 &  7167 &  7053.1  & 7167  & 0 & 95    & 3.3   & 22.1  & 25.3  \\
p07-f40-v1t2 &  7340 &  7001 &  6906.1  & 7001  & 0 & 433   & 11.0  & 75.5  & 86.5  \\
p07-f40-v2t1 &  7174 &  6881 &  6745.5  & 6881  & 0 & 317   & 8.5   & 149.0 & 157.5 \\
p07-f40-v2t2 &  7000 &  6845 &  6629.4  & 6845  & 2.34  & 43    & 20.3  & 360.1 & 380.4 \\\hline
p07-s40-v1t1 &  7860 &  7343 &  7256.3  & 7343  & 0 & 79    & 0.8   & 4.4   & 5.2   \\
p07-s40-v1t2 &  7656 &  7124 &  7091.1  & 7117  & 0 & 29    & 1.8   & 6.5   & 8.3   \\
p07-s40-v2t1 &  7466 &  7005 &  6924.5  & 7000  & 0 & 277   & 3.5   & 128.6 & 132.2 \\
p07-s40-v2t2 &  7241 &  6944 &  6762.8  & 6944  & 0 & 251   & 5.7   & 211.2 & 216.9 \\\hline
p07-t40-v1t1 &  7283 &  7013 &  6940.5  & 7012  & 0 & 206   & 10.2  & 147.8 & 158.1 \\
p07-t40-v1t2 &  7058 &  6972 &  6806.6  & 6972  & 0.35  & 66    & 76.2  & 360.5 & 436.7 \\
p07-t40-v2t1 &  6945 &  6857 &  6649.8  & 6665  & 0.04  & 19    & 68.2  & 391.1 & 459.2 \\
p07-t40-v2t2 &  6945 &  6658 &  6535.4  & 6658  & 1.64  & 6 & 26.1  & 363.3 & 389.5 \\\hline
\multicolumn{10}{l}{$^1$Optimal or best IP found by the algorithm within the time limit.}\\
\end{tabular}
\end{table}

Table  \ref{40Cust_res3} presents some details about the optimal solution: the indices of
open facilities (Loc), the number of vehicles used at each open facility (\# Veh), and the
total number of routes in each solution (\# Rt). We listed the constructed vehicle routes
and schedules (i.e., selected pairings) in the optimal (or best) solutions for all instances
in Appendix \ref{app:40customerSolutions}. In all instances 3 to 4 facilities are selected to
be open, and total number of vehicles used changes between 4 to 6. Table \ref{SolDet2}
presents some summary statistics about the solutions, such as the number of vehicles used,
the average number of routes per vehicle, the average number of customers per route and
per vehicle, and the maximum number of customers per vehicle. In an optimal or best solution
reported, a facility employs up to 3 vehicles. And each vehicle serves at least 2 routes
and may serve up to 4 routes. On average, a vehicle covers 8.3 customers and the maximum
number of customers per vehicle changes between 7 to 14.

Similar to the analysis done for 25-customer instances considering the change in the
performance of the algorithm  with vehicle capacity and time limit values,
we can analyze the results for 40-customer instances. In  7 out of 16 pairs the
computation time increases when the vehicle capacity increases, in 12 out of 16 pairs
the computation time increases when the distance limit increase. In 6 cases out of 8,
the computation time increases when both time and vehicle limits increase.
Since in one third of the 40-customer instances, the algorithm terminated with the time
limit, the solution times are not very sensitive to the changes in the capacity parameters.
On average, we calculated about \%3.8 and \%1.6 reduction in cost
with vehicle capacity increase and time limit increase, respectively. The cost
reduces about \%5.3 on average if both vehicle and time limit increases. These results
are in line with the corresponding averages for 25-customer instances.
%similar to the averages for 25-customer instances.

The solution information provided in Table
\ref{25Cust_res} and Table \ref{40Cust_res3} highlights the value of solving the integrated
problem instead of a series of problems.  Although the same number of facilities is selected
within an instance group (meaning a fixed set of customer locations and a set of candidate
facility locations), the particular facilities selected may differ depending on the vehicle
capacity and time limit values.  Which facilities are selected impacts which customers
are served by each facility and hence the sequence of the routes and the structure of the
pairings. From the solution information, we can see that, within an instance group,
increases in the vehicle capacity and the time limit reduce the total cost through changes
in facility, decreases in the number of vehicles required, changes in the underlying route
sequences.

%\begin{comment}
%Similar to the analysis done in the previous section, we summarize some similar results in
%Table \ref{EffectVCTL2} presenting the change in the performance of the algorithm
%and the solution cost with changes in vehicle capacity and the vehicle time limit.

%As expected, since the pricing problem gets more difficult, the solution time tends to increase
%with the vehicle capacity and the vehicle time limit. In addition, we expected
%the optimal cost to decrease when the vehicle capacity and/or the vehicle
%time limit increase(s). We summarize some of the results in Table \ref{EffectVCTL}. The third
%column ($L^V$, $L\uparrow$) presents the cases where time limit increases for fixed vehicle
%capacity, the fourth column ($L^V \uparrow$, $L$) presents the cases where vehicle
%capacity increases for fixed time limit, and  the fifth column ($L^V \uparrow$, $L\uparrow$)
%presents the cases both time limit and vehicle capacity increase.  Table \ref{EffectVCTL}
%shows for both 25- and 40-customer instances the number of cases in which total CPU time increased,
%the average increase in time with respect to the reference, and the average percentage
%decrease in cost with the corresponding change in vehicle capacity or time limit.
%25-customer instances seem to be more sensitive in terms of computation times to the changes in
%vehicle capacity and time limit. This is because of the shorter computational times in 25-customer
%instances and

%However, there are instances which favor the increase in time limit and/or capacity.
%An increase in the vehicle time limit can reduce the cost through a change in the routing cost
%together with a decrease in the number of vehicles and/or a change in the set of
%selected facilities. This can be seen through the solution details presented in the last three
%columns of Tables  \ref{25Cust_res} and \ref{40Cust_res2}.
%\end{comment}

\begin{table}[!htbp]
\centering
\caption{Solution details for instances with 40 customers and 5 facilities}\label{40Cust_res3}
\begin {tabular}[!h]{|c|r|r|r|c|r|r|r|}
\hline
Data file &  Loc  & \# Veh & \# Rt & Data file &  Loc & \# Veh & \# Rt \\\hline
p01-f40-v1t1 &  0,3,4 & 2,2,1   & 14    & p03-l40-v1t1  & 0,2,4 & 2,2,2 & 13    \\
p01-f40-v1t2 &  0,2,4 & 1,2,2   & 14    & p03-l40-v1t2  & 0,2,4 & 2,1,2 & 13    \\
p01-f40-v2t1 &  0,2,4 & 1,2,2   & 11    & p03-l40-v2t1  & 0,2,4 & 1,2,2 & 10    \\
p01-f40-v2t2 &  0,3,4 & 2,1,1   & 10    & p03-l40-v2t2  & 0,2,4 & 1,1,2 & 10    \\\hline
p01-l40-v1t1 &  0,2,4 & 1,2,3   & 15    & p07-f40-v1t1  & 0,2,3 & 1,2,2 & 13    \\
p01-l40-v1t2 &  0,2,4 & 1,2,2   & 15    & p07-f40-v1t2  & 0,2,4 & 1,2,2 & 13    \\
p01-l40-v2t1 &  0,2,4 & 1,2,2   & 10    & p07-f40-v2t1  & 0,2,4 & 1,2,2 & 10    \\
p01-l40-v2t2 &  0,2,4 & 1,1,2   & 10    & p07-f40-v2t2  & 1,2,4 & 1,1,2 & 10    \\\hline
p03-f40-v1t1 &  0,1,2,4 & 2,1,1,1 & 16  & p07-s40-v1t1  & 0,1,4 & 2,2,2 & 14    \\
p03-f40-v1t2 &  0,1,2,4 & 1,1,2,1 & 16  & p07-s40-v1t2  & 0,1,4 & 1,2,2 & 13    \\
p03-f40-v2t1 &  0,1,2,4 & 1,1,1,1 & 11  & p07-s40-v2t1  & 0,1,4 & 1,2,2 & 10    \\
p03-f40-v2t2 &  0,1,2,4 & 1,1,1,1 & 11  & p07-s40-v2t2  & 0,2,4 & 1,2,2 & 10    \\\hline
p03-m40-v1t1 &  0,2,4   & 2,2,2 & 15    & p07-t40-v1t1  & 0,1,4 & 1,2,2 & 13    \\
p03-m40-v1t2 &  0,2,4   & 2,2,2 & 15    & p07-t40-v1t2  & 0,2,4 & 1,2,2 & 13    \\
p03-m40-v2t1 &  0,2,4   & 2,2,2 & 11    & p07-t40-v2t1  & 0,1,4 & 1,1,2 & 10    \\
p03-m40-v2t2 &  0,1,4   & 2,1,2 & 12    & p07-t40-v2t2  & 0,2,4 & 1,1,2 & 9 \\\hline
\end{tabular}
\end{table}

\begin{table}[!htbp]
\centering
\caption{More Solution Details for 40-customer Instances}\label{SolDet2}
\begin {tabular}[!h]{|l|c|c|} \hline
%\multicolumn{1}{|c}{} &  \multicolumn{2}{c|}{40-customer instances} \\\cline{2-5}
Case                    &  Changes btw [min - max] & Avg. \\\hline
\# of vehicles per facility      & 1 - 3 &  \\
Total \# of vehicles              & 4 - 6 &  \\
Avg. \# of routes per vehicle    & 1.8 - 3.2 & 2.5 \\
Avg. \# of customers per route    & 2.5 - 4.4 & 3.4 \\
Avg. \# of customers per vehicle & 6.7 - 10 & 8.3 \\
Max \# of customers per vehicle  & 7 - 14 & 10.2 \\\hline
\end{tabular}
\end{table}

\begin{comment}
\begin{table}[b!]
\centering
\caption{Effect of Vehicle Capacity and Time Limit for 40-customer instances}\label{EffectVCTL2}
\begin {tabular}[t]{l|c|c|c}
Case & $L^V$, $L\uparrow$ & $L^V \uparrow$, $L$ & $L^V \uparrow$, $L\uparrow$ \\ \hline
%25-customer  & \# of times Total CPU increases                      & 14 / 16 & 10 / 16 & 8 / 8 \\
%\mbox{}      & Avg. increase in CPU      & 34.9 & 4.7 & 11.6 \\
%\mbox{}      & Avg. \% decrease in cost  & 3.7 & 2.2 & 5.9 \\\hline
 \#  of times Total CPU increases                       & 12 / 16 & 7 / 16 & 6 / 8 \\
Avg. increase in CPU      & 2.7 & 4.8 & 8.7 \\
 Avg. \% decrease in cost  & 1.6 & 3.8 & 5.3 \\\hline
\end{tabular}
\end{table}
\end{comment}


\subsection{Performance of Branch and Price Algorithm for Instances from Literature}
\label{compareLin}
The only LRSP instances available in the literature are those of \citet{Lin};
%we present our results for those instances in Section~\ref{compareLin}.
%For further testing, we generated several sets of larger instances.
%We describe the instances in Section~\ref{instances} and then
%present the results in Sections \ref{alternativeForm} and \ref{bp-results}.
%\subsection{Instances from \citet{Lin}}
\citet{Lin} present six instances: a base instance that includes
27 customers and four facility locations in Hong Kong plus five instances
generated from the original by considering three subsets of 10 customers and two subsets
of 12 customers while keeping the same four facility locations. We used only
the four instances available to the reader to test our branch-and-price algorithm,
in particular, our E-BP algorithm.

\citet{Lin} solve these instances with a branch-and-bound algorithm
and with their heuristic algorithm.  Within a time limit of 10,000 seconds, their
branch-and-bound algorithm is able to prove optimality only for the 10 customer
instances and is not able to find a feasible solution for the original 27 customer problem.
The heuristic quickly finds solutions for all instances.
Table \ref{CompLin} compares their reported results
to results obtained with E-BP. As shown in Table \ref{CompLin},
our algorithm quickly provides optimal solutions for all four instances.  (The * indicates
the best solution found in 10,000 CPU seconds on a Pentium III machine.)

%[This would be better as one table.  Also, you need to explain or change the 309,816 objective
%function for the first instance.]
\begin{table}[h!tbp]
  \begin{center}
  \caption{Comparison of \citet{Lin} results and branch-and-price results} \label{CompLin}
  \begin {tabular}[h]{|c|c|c|c|c|c|c|c|} \hline
\multicolumn{2}{|c}{Instance} & \multicolumn{2}{|c}{Lin B \& B} & \multicolumn{2}{|c}{Lin Heuristic} & \multicolumn{2}{|c|}{Ak\c{c}a E-BP} \\ \hline
\# Cust & \# Fac & Objective & CPU (s) & Objective & CPU(s) & Objective & CPU(s) \\ \hline
10 & 4 & 309,817   &    1155    &  309,817  &   0.44    & 309,817   &  0.66\\
10 & 4 & 309,808   &    982 & 309,808   &   0.49    & 309,808       & 0.10  \\
12 & 4 & 312,036*  & $>$ 10,000 &  312,036  &   0.82    & 312,036       & 0.05\\
27 & 4 & -         &      -     & 625,752.5 &     6     & 625,750.2 & 57.83 \\\hline
\end{tabular}
\end{center}
\end{table}


\begin{comment}
%\section{Branch-and-Price Results}
%\label{bp-results}
In our main experiments, we evaluated the performance of the branch-and-price
algorithm on the 25- and 40-customer instances.
For the 25-customer instances, we used E-BP and applied a time limit of 6 CPU hours. Table
\ref{25Cust_res} shows the results for instances 1 through 8 with 25 customers.
For each instance, Table \ref{25Cust_res} reports the instance number, the vehicle
capacity ($L^V$), the time limit ($L^T$), the upper bound obtained from I-Heur (I UB),
the LP relaxation value at the root node (LP), the optimal objective value (IP),
the integrality gap at the end of the running time (Gap \%), the number of evaluated
nodes (Nds), and the total CPU time in minutes (CPU). The last three columns present
some details about the optimal solution: the indices of open facilities
(Loc), the number of vehicles used at each open facility (\# Veh), and the total number
of routes in each solution (\# Rt). Overall, the algorithm handled the 25-customer
instances well; 28 of the
32 instances were solved to optimality in less than 25 minutes with the majority (21)
solved in 10 minutes or less. For only one instance did the algorithm terminate due to the
time limit before proving optimality.

For the 40-customer instances, we used E-BP-2S with a time limit of 2 CPU hours for
stage one and 6 CPU hours for stage two and reported the results in Table \ref{40Cust_res2}.
The labels in Table \ref{40Cust_res2} are same as the ones in Table \ref{25Cust_res}
but with three additional columns: the fifth column (H-BP UB) contains
the upper bound obtained at the end of H-BP, the tenth and eleventh columns contain
the CPU times spent in H-BP and E-BP, respectively. As expected,
the solution times in the 40-customer instances were larger than the 25-customer instances.
Even though the algorithm terminated with a nonzero gap in some cases,
the algorithm was able to prove optimality most of the time.
For other instances, good solutions with small integrality gaps were obtained
within the time limit.  The gap was at most 0.35\% for four of the instances and between
1.26\% and 4.18\% for the other seven instances.

The initial feasible solutions provided by I-Heur were good for the 25-customer
instances with an average deviation of 4.4\% from optimality. Similarly,
for the 40-customer instances, the average deviation of the upper bound obtained
from I-Heur from the optimal or best integer solution was 4.5\%.
However, for the 40-customer instances, the overall success of the algorithm can be
attributed to the effectiveness of H-BP in finding good upper bounds. In all but one
instance, the H-BP upper bound was better than the I-Heur upper bound, and for 23 instances
out of 32, the H-BP upper bound value equaled that of the optimal (or best integer) solution.

The LP relaxation bounds were quite strong, as expected with the addition of valid inequalities
at the root node. The average deviation of the LP relaxation bound from the optimality
was 2.6\% in 25-customer instances and 2.2\% in 40-customer instances, which are in line with
the results of the smaller instances (instances with 20 customers in Section \ref{alternativeForm}).

In addition to the information in Table  \ref{25Cust_res} and Table \ref{40Cust_res2},
Table \ref{SolDet}
presents some summary statistics about the solutions, such as the number of vehicles used,
the average number of routes
per vehicle, the average number of customers per route and per vehicle, and the maximum number
of customers per vehicle.

The solution information provided in Table \ref{25Cust_res} and Table \ref{40Cust_res2} highlights
the value of solving the integrated problem instead of a series of problems.  Although the same number
of facilities is selected within an instance group (meaning a fixed set of customer locations and
a set of candidate facility locations), the particular facilities selected may differ depending
on the vehicle capacity and time limit values.  Which facilities are selected impacts which customers
are served by each facility and hence the sequence of the routes and the structure of the pairings.
From the solution information, we can see that, within an instance group, increases in the vehicle
capacity and the time limit reduce the total cost through changes in facility, decreases in the
number of vehicles required, changes in the underlying route sequences.


%\begin{comment}
As expected, since the pricing problem gets more difficult, the solution time tends to increase
with the vehicle capacity and the vehicle time limit. In addition, we expected
the optimal cost to decrease when the vehicle capacity and/or the vehicle
time limit increase(s). We summarize some of the results in Table \ref{EffectVCTL}. The third
column ($L^V$, $L\uparrow$) presents the cases where time limit increases for fixed vehicle
capacity, the fourth column ($L^V \uparrow$, $L$) presents the cases where vehicle
capacity increases for fixed time limit, and  the fifth column ($L^V \uparrow$, $L\uparrow$)
presents the cases both time limit and vehicle capacity increase.  Table \ref{EffectVCTL}
shows for both 25- and 40-customer instances the number of cases in which total CPU time increased,
the average increase in time with respect to the reference, and the average percentage
decrease in cost with the corresponding change in vehicle capacity or time limit.
25-customer instances seem to be more sensitive in terms of computation times to the changes in
vehicle capacity and time limit. This is because of the shorter computational times in 25-customer
instances and in one third of the 40-customer instances, the algorithm terminated with the time limit.
However, there are instances which favor the increase in time limit and/or capacity.
An increase in the vehicle time limit can reduce the cost through a change in the routing cost
together with a decrease in the number of vehicles and/or a change in the set of
selected facilities. This can be seen through the solution details presented in the last three
columns of Tables  \ref{25Cust_res} and \ref{40Cust_res2}.
\end{comment}

\begin{comment}
\begin{table}[h!tbp]
\centering
\caption{Results for instances with 25 customers and 5 facilities}\label{25Cust_res}
\begin {tabular}[t]{|c|c|c|c|c|c|c|r|r|c|c|c|}
\hline
\# & $L^V$ & $L^T$ & I UB & LP & IP$^1$ & Gap \% & Nds & CPU (m) & Loc & \# Veh & \# Rt \\ \hline
1 & 50  & 140   & 4795  & 4504.4 & 4699 & 0  & 31   & 0.9   & 0,4   & 2,2   & 10    \\
1 & 50  & 160   & 4795  & 4425.8 & 4511 & 0 & 197   & 280.5 & 0,4   & 1,2   & 10    \\
1 & 70  & 140   & 4722  & 4366.1 & 4425 & 0 & 97    & 24.4  & 0,4   & 1,2   & 7 \\
1 & 70  & 160   & 4722  & 4295.4 & 4425 & 0 & 71    & 20.4  & 0,4   & 1,2   & 7 \\\hline
2 & 50  & 140   & 5193  & 4849.8 & 4899 & 0 & 22    & 0.7   & 0,1   & 2,2   & 8 \\
2 & 50  & 160   & 4968  & 4725.4 & 4815 & 0 & 121   & 3.3   & 0,4   & 2,2   & 8 \\
2 & 70  & 140   & 4809  & 4669.0 & 4759 & 0 & 157   & 5.8   & 0,4   & 2,2   & 6 \\
2 & 70  & 160   & 4762  & 4537.1 & 4602 & 0 & 39    & 8.6   & 0,1   & 1,2   & 7 \\\hline
3 & 60  & 140   & 5281  & 4810.0 & 5062 & 0 & 268   & 2.7   & 0,2   & 3,2   & 10    \\
3 & 60  & 160   & 5139  & 4701.1 & 4837 & 0 & 75    & 2.2   & 0,2   & 2,2   & 10    \\
3 & 80  & 140   & 5048  & 4606.1 & 4758 & 0 & 579   & 5.1   & 0,4   & 2,2   & 8 \\
3 & 80  & 160   & 4831  & 4502.5 & 4742 & 0 & 49    & 6.7   & 0,4   & 2,2   & 8 \\\hline
4 & 60  & 140   & 4995  & 4697.4 & 4785 & 0 & 567   & 5.7   & 0,4   & 2,2   & 10    \\
4 & 60  & 160   & 4869  & 4578.2 & 4785 & 0.1 & 3790    & 360.0 & 0,4   & 2,2   & 10    \\
4 & 80  & 140   & 4992  & 4485.8 & 4689 & 0 & 59    & 3.6   & 0,4   & 2,2   & 7 \\
4 & 80  & 160   & 4767  & 4396.1 & 4479 & 0 & 9 & 10.6  & 0,4   & 1,2   & 7 \\\hline
5 & 60  & 140   & 5186  & 4795.3 & 4842 & 0 & 17    & 0.6   & 0,2   & 2,2   & 9 \\
5 & 60  & 160   & 4871  & 4670.2 & 4803 & 0 & 35    & 1.9   & 0,4   & 2,2   & 7 \\
5 & 80  & 140   & 4805  & 4620.9 & 4741 & 0 & 43    & 2.8   & 0,4   & 2,2   & 6 \\
5 & 80  & 160   & 4805  & 4507.6 & 4655 & 0 & 45    & 11.5  & 2,4   & 1,2   & 7 \\\hline
6 & 50  & 140   & 4870  & 4642.6 & 4761 & 0 & 57    & 1.6   & 0,2   & 1,3   & 7 \\
6 & 50  & 160   & 4823  & 4544.5 & 4685 & 0 & 71    & 130.6 & 2,4   & 2,1   & 9 \\
6 & 70  & 140   & 4721  & 4485.6 & 4703 & 0 & 263   & 9.5   & 0,2   & 2,2   & 6 \\
6 & 70  & 160   & 4721  & 4388.3 & 4482 & 0 & 181   & 23.9  & 0,2   & 1,2   & 6 \\\hline
7 & 50  & 140   & 5412  & 5000.1 & 5098 & 0 & 441   & 1.1   & 0,4   & 2,3   & 9 \\
7 & 50  & 160   & 5238  & 4852.2 & 4889 & 0 & 161   & 3.2   & 0,4   & 2,2   & 9 \\
7 & 70  & 140   & 5159  & 4739.4 & 4787 & 0 & 21    & 0.5   & 0,4   & 2,2   & 7 \\
7 & 70  & 160   & 4934  & 4609.8 & 4781 & 0 & 105   & 4.0   & 0,4   & 2,2   & 8 \\\hline
8 & 50  & 140   & 5134  & 4836.5 & 4898 & 0 & 39    & 1.3   & 0,1   & 2,2   & 8 \\
8 & 50  & 160   & 5162  & 4703.1 & 4886 & 0 & 2475  & 91.8  & 0,2   & 2,2   & 8 \\
8 & 70  & 140   & 5025  & 4646.2 & 4767 & 0 & 781   & 16.5  & 0,2   & 2,2   & 6 \\
8 & 70  & 160   & 4769  & 4513.1 & 4592 & 0 & 149   & 20.7  & 0,1   & 1,2   & 6 \\\hline
\multicolumn{11}{l}{$^1$Optimal or best IP found by the algorithm within the time limit.}\\
\end{tabular}
\end{table}
\begin{table}[h!tbp]
\centering
\caption{Results for instances with 40 customers and 5 facilities}\label{40Cust_res2}
\begin {tabular}[t]{|c|c|c|c|c|c|c|r|r|r|r|r|c|c|c|}
\hline
  &       &       & I  & H-BP &    &        &    Gap     &       & \multicolumn{3}{|c|}{CPU(m)} &  & \#  & \#  \\ \cline{10-12}
\#& $L^V$ & $L^T$ & UB & UB   & LP & IP$^1$ & \% \ \ & Nds & H-BP & E-BP & Total & Loc & Veh &  Rt \\ \hline
9 & 50  & 140   & 7170  & 7131  & 6976.0 & 7131 & 0 & 807   & 53.2  & 150.0 & 203.2 & 0,3,4 & 2,2,1 & 14    \\
9 & 50  & 160   & 7170  & 7170  & 6826.2 & 6999 & 0.32  & 1065  & 120.0 & 360.0 & 480.0 & 0,2,4 & 1,2,2 & 14    \\
9 & 70  & 140   & 7064  & 6866  & 6700.2 & 6866 & 0 & 203   & 4.9   & 113.0 & 117.9 & 0,2,4 & 1,2,2 & 11    \\
9 & 70  & 160   & 6984  & 6821  & 6581.0 & 6821 & 1.26  & 35    & 35.4  & 360.7 & 396.1 & 0,3,4 & 2,1,1 & 10    \\\hline
10  & 50 & 140  & 7449  & 7183  & 6899.7 & 7183 & 3.22  & 573   & 120.1 & 360.4 & 480.5 & 0,2,4 & 1,2,3 & 15    \\
10 & 50  & 160  & 7236  & 6934  & 6757.4 & 6934 & 0 & 43    & 5.6   & 48.0  & 53.6  & 0,2,4 & 1,2,2 & 15    \\
10 & 70 & 140   & 7166  & 6823  & 6616.7 & 6823 & 0 & 129   & 2.4   & 103.2 & 105.6 & 0,2,4 & 1,2,2 & 10    \\
10 & 70 & 160   & 7166  & 6633  & 6513.6 & 6612 & 0 & 25    & 8.4   & 187.6 & 196.0 & 0,2,4 & 1,1,2 & 10    \\\hline
11 & 60 & 140   & 9253  & 8780  & 8697.4 & 8780 & 0 & 97    & 9.6   & 97.8  & 107.4 & 0,1,2,4 & 2,1,1,1 &   16  \\
11 & 60 & 160   & 9028  & 8753  & 8578.6 & 8753 & 0 & 41    & 28.0  & 154.6 & 182.6 & 0,1,2,4 & 1,1,2,1 &   16  \\
11 & 80 & 140   & 8943  & 8663  & 8439.0 & 8440 & 0 & 3 & 7.8   & 8.8   & 16.6  & 0,1,2,4 & 1,1,1,1 &   11  \\
11 & 80 & 160   & 8943  & 8438  & 8390.4 & 8438 & 0 & 9 & 11.5  & 138.2 & 149.7 & 0,1,2,4 & 1,1,1,1 &   11  \\\hline
12 & 60 & 140   & 7827  & 7740  & 7050.9 & 7253 & 0 & 2524  & 120.0 & 360.1 & 480.1 & 0,2,4 & 2,2,2 & 15    \\
12 & 60 & 160   & 7544  & 7236  & 6907.7 & 7236 & 2.95  & 24    & 84.0  & 360.4 & 444.4 & 0,2,4 & 2,2,2 & 15    \\
12 & 80 & 140   & 7479  & 7132  & 6793.9 & 7132 & 2.79  & 27    & 120.0 & 360.1 & 480.1 & 0,2,4 & 2,2,2 & 11    \\
12 & 80 & 160   & 7430  & 6947  & 6655.5 & 6947 & 4.18  & 4 & 37.2  & 360.3 & 397.5 & 0,1,4 & 2,1,2 & 12    \\\hline
13 & 60 & 140   & 7386  & 7167  & 6920.4 & 7167 & 0 & 183   & 1.5   & 55.6  & 57.1  & 0,2,4 & 2,2,2 & 13    \\
13 & 60 & 160   & 7212  & 6950  & 6771.6 & 6950 & 0.05  & 2121  & 120.0 & 360.0 & 480.0 & 0,2,4 & 2,1,2 & 13    \\
13 & 80 & 140   & 7121  & 6849  & 6662.9 & 6849 & 0 & 159   & 3.9   & 113.8 & 117.7 & 0,2,4 & 1,2,2 & 10    \\
13 & 80 & 160   & 7018  & 6846  & 6544.7 & 6639 & 0 & 53    & 39.0  & 290.9 & 329.9 & 0,2,4 & 1,1,2 & 10    \\\hline
14 & 50 & 140   & 7424  & 7167  & 7053.1 & 7167 & 0 & 95    & 3.3   & 22.1  & 25.3  & 0,2,3 & 1,2,2 & 13    \\
14 & 50 & 160   & 7340  & 7001  & 6906.1 & 7001 & 0 & 433   &11.0   & 75.5  & 86.5  & 0,2,4 & 1,2,2 & 13    \\
14 & 70 & 140   & 7174  & 6881  & 6745.5 & 6881 & 0 & 317   & 8.5   & 149.0 & 157.5 & 0,2,4 & 1,2,2 & 10    \\
14 & 70 & 160   & 7000  & 6845  & 6629.4 & 6845 & 2.34  & 43    & 20.3  & 360.1 & 380.4 & 1,2,4 & 1,1,2 & 10    \\\hline
15 & 50 & 140   & 7860  & 7343  & 7256.3 & 7343 & 0 & 79    & 0.8   & 4.4   & 5.2   & 0,1,4 & 2,2,2 & 14    \\
15 & 50 & 160   & 7656  & 7124  & 7091.1 & 7117 & 0 & 29    & 1.8   & 6.5   & 8.3   & 0,1,4 & 1,2,2 & 13    \\
15 & 70 & 140   & 7466  & 7005  & 6924.5 & 7000 & 0 & 277   & 3.5   & 128.6 & 132.2 & 0,1,4 & 1,2,2 & 10    \\
15 & 70 & 160   & 7241  & 6944  & 6762.8 & 6944 & 0 & 251   & 5.7   & 211.2 & 216.9 & 0,2,4 & 1,2,2 & 10    \\\hline
16 & 50 & 140   & 7283  & 7013  & 6940.5 & 7012 & 0 & 206   & 10.2  & 147.8 & 158.1 & 0,1,4 & 1,2,2 & 13    \\
16 & 50 & 160   & 7058  & 6972  & 6806.6 & 6972 & 0.35  & 66    & 76.2  & 360.5 & 436.7 & 0,2,4 & 1,2,2 & 13    \\
16 & 70 & 140   & 6945  & 6857  & 6649.8 & 6665 & 0.04  & 19    & 68.2  & 391.1 & 459.2 & 0,1,4 & 1,1,2 & 10    \\
16 & 70 & 160   & 6945  & 6658  & 6535.4 & 6658 & 1.64 & 6  & 26.1  & 363.3 & 389.5 & 0,2,4 & 1,1,2 & 9 \\\hline
\multicolumn{15}{l}{$^1$Optimal or best IP found by the algorithm within the time limit.}\\
\end{tabular}
\end{table}
%\end{landscape}
\begin{table}[b!]
\centering
\caption{Effect of Vehicle Capacity and Time Limit}\label{EffectVCTL2}
\begin {tabular}[t]{l|l|c|c|c}
\mbox{} & \mbox{} & $L^V$, $L\uparrow$ & $L^V \uparrow$, $L$ & $L^V \uparrow$, $L\uparrow$ \\ \hline
%25-customer  & \# of times Total CPU increases                      & 14 / 16 & 10 / 16 & 8 / 8 \\
%\mbox{}      & Avg. increase in CPU      & 34.9 & 4.7 & 11.6 \\
%\mbox{}      & Avg. \% decrease in cost  & 3.7 & 2.2 & 5.9 \\\hline
40-customer  & \#  of times Total CPU increases                       & 12 / 16 & 7 / 16 & 6 / 8 \\
\mbox{}      & Avg. increase in CPU      & 2.7 & 4.8 & 8.7 \\
\mbox{}      & Avg. \% decrease in cost  & 1.6 & 3.8 & 5.3 \\\hline
\end{tabular}
\end{table}
\begin{table}[h!tbp]
\centering
\caption{More Solution Details}\label{SolDet}
\begin {tabular}[t]{|l|c|c||c|c|} \hline
\multicolumn{1}{|c}{} & \multicolumn{2}{|c||}{25-customer instances} & \multicolumn{2}{c|}{40-customer instances} \\\cline{2-5}
                          & min - max & Avg. &  min - max & Avg. \\\hline
\# of vehicles per facility      &  1 - 3 &  & 1 - 3 &  \\
Total \# of vehicles             &  3 - 5  &  & 4 - 6 &  \\
Avg. \# of routes per vehicle    &  1.5 - 3.3 & 2.1 & 1.8 - 3.2 & 2.5 \\
Avg. \# of customers per route   &  2.5 - 4.2 & 3.3 & 2.5 - 4.4 & 3.4 \\
Avg. \# of customers per vehicle &  5 - 8.3   & 6.9 & 6.7 - 10 & 8.3 \\
Max \# of customers per vehicle  &  7 - 11    & 8.5 & 7 - 14 & 10.2 \\\hline
\end{tabular}
\end{table}
\end{comment}



%\zaupdate{\subsection{sensitivity to costs: Facility fixed cost vehicle fixed cost}}
%\zaupdate{ meeting reports/comresultsfeb807}

%\zaupdate{\subsection{Performance of Branching Rule Selection}}
%Order of branching rules and selection specific rule }}
%\zaupdate{ meeting reports/Dec13-06}


\section{Computational Experiments For Chapter 3}
The goal of the computational experiments in this section was
to test the performance of the branch and price algorithm extended with
some of the cutting planes described in Chapter 3. We empirically show the
effect of each implemented type of cut. We implemented lifted cover
inequalities, generalized assignment polytope cuts and disjunctive cuts.

\subsection{Lifted Cover Inequalities and GAP Cuts}
\label{comp:CoverGAP}
In this section, we evaluated the effect of lifted cover inequalities
(inequalities in the form of \ref{inq-cover2c}) and the GAP cuts
(inequalities in the form of \ref{gap-valid2}) which are valid for
the projections (knapsack polytope and generalized assignment polytope)
of SPP-LRSP. The details of these inequalities were described in Section
\ref{pro-spp}.

For the experiments in this Section, we used eight 25-customer instances
with largest vehicle capacity and time limit (x-v2t2 instances).
We used 1-step branch and price algorithm with 2p-ESPPRC-CS(12),
2p-ESPPRC-LL(50),  2p-ESPPRC-LL(100), 2p-ESPPRC-LL(200) and
2p-ESPPRC pricing algorithms. We ran cut generation at each node if it
its depth is less than 9. At each node the steps of the algorithm extended
with cut generation is as follows:
\begin{itemize}
\item Call COLGEN.
\item If the LP is fractional and depth $\leq$ 8, call CUTGEN, otherwise STOP.
\item If a cut is generated and added to the model, call COLGEN,
otherwise STOP.
 \end{itemize}
Here, we refer cut generation function as CUTGEN. Depending on the type of
the cut we wanted, CUTGEN procedure searches for either lifted cover inequalities
or GAP cuts, or both. We refer these procedures as CUTGEN(KNAP),
CUTGEN(GAP) and CUTGEN(KNAP,GAP). As explained in Section \ref{pro-spp},
at each node, CUTGEN procedure projects the LP solution from $(t,z)$ space to
$(t,v)$ space and uses separation algorithms given in the Section to generate
required type of cuts. If the new cuts are generated, they are
projected back to the original space as in (\ref{inq-cover2c}) and (\ref{gap-valid2}).
At each node, we called CUTGEN only once and added
generated cuts all at once before we do column generation again to find
the optimal LP relaxation solution with cuts at the evaluated node.
In addition, these cuts are global cuts.

Table \ref{knpgap25} compares the performance of the algorithms with
CUTGEN(KNAP) and CUTGEN(GAP) with the base algorithm without any cuts.
Table \ref{knpgap25} reports the number of evaluated nodes (Column Nds)
and the number of cuts generated in each algorithm. In addition, Table gives
the ratio of the time of the algorithm with CUTGEN to the time of the algorithm
without cuts (column Time W/WO).

Table \ref{knpgap25} shows that CUTGEN procedure slightly effects the
performance of the algorithm.  In all of the instances except one case,
if the cut generation is successful the number of evaluated nodes
either stayed same or decreased. On average, the computational time
of the algorithms with CUTGEN is almost same to the base algorithm.
On each instance, the time in the algorithms with CUTGEN is either slightly
lower or slightly higher than the time in the base algorithm.

\begin{table}[h!tbp]
\centering
\caption{Cover Inequalities and GAP Cuts for 25-customer Instances}\label{knpgap25}
\begin {tabular}[t]{|c|c|c|c|c|c|c|c|}
\hline
\mbox{} & Base Alg. & \multicolumn{3}{|c|}{Alg. with GAP Cuts} & \multicolumn{3}{|c|}{Alg. with COVER Cuts} \\\cline{2-8}
Data file & Nds & Nds & Time W/WO & \# cuts & Nds & Time W/WO & \# cuts \\\hline
p01-f25-v2t2 &  29 &    29 &    1       & 0 &   41 &    0.97    & 3 \\
p01-l25-v2t2 &  33 &    33 &    0.96    & 1 &   33 &    0.99    & 1 \\
p03-f25-v2t2 &  44 &    40 &    1.03    & 5 &   36 &    1.14    & 14    \\
p03-m25-v2t2 &  9 & 9 & 1.07    & 0 &   9  &    1.02    & 0 \\
p03-l25-v2t2 &  37 &    37 &    1.04    & 0 &   37 &    1.06    & 4 \\
p07-f25-v2t2 &  143 &   143 &   1.06    & 0 &   143 &   1.02    & 0 \\
p07-s25-v2t2 &  85 &    83 &    0.97    & 1 &   81 &    1   & 5 \\
p07-t25-v2t2 &  91 &    81 &    0.94    & 2 &   69 &    0.9 & 7 \\\hline
avg & \mbox{} & \mbox{} &   1.01    & \mbox{} & \mbox{} & 1.01 & \mbox{} \\\hline
\end{tabular}
\end{table}

Table \ref{knpgap25-2} presents the results of the algorithm with CUTGEN(KNAP,GAP)
and compares it with the base algorithm. The labels in the Table are same as
Table \ref{knpgap25}. The results are parallel to the algorithms  with CUTGEN(KNAP)
and CUTGEN(GAP). Even though the LP relaxation bound improves slightly
with the generated cuts at the relevant nodes, the CUTGEN procedure does not have a
significant contribution to the performance of the overall algorithm in this experiment.
%One factor in this
\begin{table}[h!tbp]
\centering
\caption{Cover Inequalities and GAP Cuts for 25-customer Instances}\label{knpgap25-2}
\begin {tabular}[t]{|c|c|c|c|c|c|}
\hline
\mbox{} & Base Alg. & \multicolumn{4}{|c|}{Alg. with COVER + GAP Cuts} \\\cline{2-6}
Data file & Nds & Nds & Time W/WO & \# gap cuts & \# cover cuts \\\hline
p01-f25-v2t2 &  29  & 41    & 0.9   & 0 & 3 \\
p01-l25-v2t2 &  33  & 33    & 0.99  & 1 & 1 \\
p03-f25-v2t2 &  44  & 36    & 1.14  & 0 & 14    \\
p03-m25-v2t2 &  9   & 9 & 1.03  & 0 & 0 \\
p03-l25-v2t2 &  37  & 37    & 1.05  & 0 & 4 \\
p07-f25-v2t2 &  143 & 143   & 1.01  & 0 & 0 \\
p07-s25-v2t2 &  85  & 81    & 0.99  & 0 & 5 \\
p07-t25-v2t2 &  91  & 69    & 0.9   & 0 & 7 \\\hline
avg & \mbox{}   & \mbox{}   & 1 & \mbox{}   & \mbox{} \\\hline
\end{tabular}
\end{table}

\subsection{Disjunctive Cuts}
\label{sec:disjcutres}

In this Section, we present some computational results for four
implementations of disjunctive cutting plane generation algorithm
described in Section \ref{disjcuts-sec}. We refer the cut generation
procedures for each of four implementations as DISJCUTGEN(1),
DISJCUTGEN(2), DISJCUTGEN(3), and DISJCUTGEN(4).

In this set of experiments, we used the same 1 step branch and price algorithm
design as in Section \ref{comp:CoverGAP}, except we used DISJCUTGEN procedures
instead of CUTGEN procedures. Again, at each node, we called the cut generation
procedure at most once up to depth 8 in the tree. Another difference was that we
added the generated cuts as local cuts for the corresponding node.
In DISJCUTGEN(2), DISJCUTGEN(3) and DISJCUTGEN(4) (second, third and fourth
implementations), at each node we searched one set of disjunctive inequality and
at most one disjunctive cut. In DISJCUTGEN(1) (first disjunctive cutting plane
implementation described in Section \ref{sec:djc-imp1}), at each node when we called
the cut generation function, we searched at most two sets of disjunctions in the form of
(\ref{imp1-disj1}) and (\ref{imp1-disj2}) and searched for at most four
cuts.

We compared the performance of each implementation of DISJCUTGEN using eight
25-customer instances. Table \ref{tab:disj-imp12} evaluates the effect of
using DISJCUTGEN(1) and  DISJCUTGEN(2), and Table \ref{tab:disj-imp34}
evaluates the effect of using DISJCUTGEN(3) and  DISJCUTGEN(4). In Tables,
base algorithm refers to the algorithm without any cutting planes.
Both Tables \ref{tab:disj-imp12} and  \ref{tab:disj-imp34} present
the number of evaluated nodes (Nds) in each algorithm, fraction
of time of each algorithm with disjunctive cuts to the base algorithm
(Time W/WO), and the number of total cuts (\# cuts) generated in each
algorithm. In all four implementations, in most of the instances, the
number of evaluated nodes stays same or increases compared to the base algorithm.
On average,  the algorithm time increases 8\% with DISJCUTGEN(1) and
25\% with DISJCUTGEN(2). DISJCUTGEN(3) increases the total algorithm time
about 3 times while DISJCUTGEN(4) increases the algorithm time significantly,
about 13 times.


To evaluate the performance reduction with DISJCUTGEN procedures further,
in Table \ref{tab:disj-time}, for each implementation and for each instance,
we present percentage of total algorithm time spent to generate cuts
(DISJCUTGEN time /Total), and  average cut generation time per cut
(DISJCUTGEN time per cut (s)). Time spent in DISJCUTGEN procedure is
mostly the time to solve the CGLPs. Since the CGLP in implementation
1 is very small compared to other implementations, the time spent
in DISJCUTGEN is negligible. CGLP in implementation 2 is larger than
implementation 1, but much smaller than the CGLPs in implementations 3 and 4
with respect to number of constraints and variables.
The percentage of total time spent in DISJCUTGEN procedure increases significantly
in implementation 3 and 4. The number of variables and main set of constraints
in CGLPs for implementations 3 and 4 are same. However, in implementation 4,
we also added some restrictions for the coefficients of variables
in the generated cut, this makes the CGLP more difficult to solve and results
more significant increase in time for DISJCUTGEN.
These results can also be seen from the reported DISJCUTGEN times per cut.

\begin{table}[!htbp]
\centering
\caption{Disjunctive Cutting Planes: Implementation 1 and 2 }\label{tab:disj-imp12}
\begin {tabular}[!h]{|c|c|c|c|c|c|c|c|}
\hline
\mbox{} & Base Alg. & \multicolumn{3}{|c|}{Alg. with DISJCUTGEN(1)} & \multicolumn{3}{|c|}{Alg. with DISJCUTGEN(2)} \\\cline{2-8}
%\mbox{} & \mbox{} & \mbox{} & Time & Cutgen/ & \mbox{} & \mbox{} & Time & Cutgen/ &   \mbox{} \\
Data file & Nds  & Nds & Time W/WO & \# cuts & Nds & Time W/WO & \# cuts \\\hline
p01-f25-v2t2 &  29  & 69    & 1.01  & 7 & 40    & 1.25  & 10    \\
p01-l25-v2t2 &  33  & 35    & 1.15  & 9 & 33    & 1.17  & 10    \\
p03-f25-v2t2 &  44  & - & - & 18    & 50    & 1.59  & 27    \\
p03-m25-v2t2 &  9   & 9 & 1.27  & 6 & 11    & 1.29  & 6 \\
p03-l25-v2t2 &  37  & 39    & 1.09  & 10    & 39    & 1.16  & 15    \\
p07-f25-v2t2 &  143 & 143   & 0.98  & 0 & 119   & 1.18  & 11    \\
p07-s25-v2t2 &  85  & 96    & 1.08  & 30    & 98    & 1.35  & 26    \\
p07-t25-v2t2 &  91  & 85    & 0.96  & 16    & 93    & 0.99  & 12    \\\hline
avg & \mbox{}   & \mbox{} & 1.08 & \mbox{} & \mbox{}    & 1.25  & \mbox{}\\\hline
\multicolumn{8}{l}{- terminates since number of non zero coefficients for constraints $\geq$ 76000.}
\end{tabular}
\end{table}

\begin{table}[!htbp]
\centering
\caption{Disjunctive Cutting Planes: Implementation 3 and 4 }\label{tab:disj-imp34}
\begin {tabular}[!h]{|c|c|c|c|c|c|c|c|}
\hline
\mbox{} & Base Alg. & \multicolumn{3}{|c|}{Alg. with DISJCUTGEN(3)} & \multicolumn{3}{|c|}{Alg. with DISJCUTGEN(4)} \\\cline{2-8}
Data file & Nds  & Nds & Time W/WO & \# cuts & Nds & Time W/WO & \# cuts \\\hline
p01-f25-v2t2 &  29  & 71    & 2.31  & 34    & 41    & 7.06  & 13 \\
p01-l25-v2t2 &  33  & 37    & 2.18  & 28    & 31    & 8.38  & 11 \\
p03-f25-v2t2 &  44  & 591   & 6.26  & 41    & 44    & 29.26 & 32 \\
p03-m25-v2t2 &  9   & 9 & 1.78  & 7 & 7 & 3.83  & 5 \\
p03-l25-v2t2 &  37  & 43    & 1.89  & 23    & 45    & 10.69 & 19 \\
p07-f25-v2t2 &  143 & 53    & 1.66  & 28    & 157   & 8 & 12 \\
p07-s25-v2t2 &  85  & 213   & 4.48  & 66    & 95    & 31.38 & 37 \\
p07-t25-v2t2 &  91  & 52    & 1.85  & 37    & 83    & 6.08  & 19 \\\hline
avg & \mbox{} & \mbox{}& 2.8 & \mbox{}& \mbox{} & 13.08 & \mbox{} \\\hline
\end{tabular}
\end{table}

\begin{table}[!htbp]
\centering
\caption{Time in DISJCUTGEN procedure for implementations 1, 2, 3 and 4 }\label{tab:disj-time}
\begin {tabular}[!h]{|c|c|c|c|c|c|c|c|c|}
\hline
\mbox{} & \multicolumn{4}{|c|}{DISJCUTGEN time /Total } & \multicolumn{4}{|c|}{DISJCUTGEN time per cut (s)} \\\cline{2-9}
Data file & Imp.1  & Imp.2  & Imp.3  & Imp.4 & Imp.1  & Imp.2  & Imp.3  & Imp.4 \\\hline
p01-f25-v2t2 & 0 & 0.02 & 0.2   & 0.6   & 0.1   & 1.46  & 7.89  & 191.41 \\
p01-l25-v2t2 & 0 & 0.02 & 0.31  & 0.8   & 0.05  & 0.88  & 10.62 & 264.22 \\
p03-f25-v2t2 & 0 & 0.05 & 0.47  & 0.91  & 0 & 0.85  & 20.71 & 243.07 \\
p03-m25-v2t2 & 0 & 0.01 & 0.09  & 0.37  & 0.02  & 0.52  & 10.24 & 126.99 \\
p03-l25-v2t2 & 0 & 0.02 & 0.33  & 0.81  & 0.05  & 0.68  & 12.49 & 211.58 \\
p07-f25-v2t2 & 0 & 0.03 & 0.25  & 0.75  &  -    & 1.98  & 9.9   & 329.24 \\
p07-s25-v2t2 & 0 & 0.09 & 0.66  & 0.94  & 0.04  & 1.03  & 9.94  & 176.54 \\
p07-t25-v2t2 & 0 & 0.04 & 0.31  & 0.7   & 0.12  & 1.92  & 9.96  & 142.53 \\\hline
avg & 0 & 0.03 & 0.33 & 0.74 & 0.05 & 1.16  & 11.47 & 210.7 \\\hline
\end{tabular}
\end{table}

The presented results showed that the DISJCUTGEN procedure reduces the algorithm
performance. However, even though the time increases in other parts of the main
algorithm, are these cuts helping somehow by improving the LP relaxation
bound at the related nodes? To be able see the result of this question, we
looked at the optimal LP solutions before and after adding the cuts in each
implementation. \zaupdate{maybe prepare a table?}

In {\em implementation 3}, almost none of the cuts violates the current
fractional LP solution of RMP ($\bar{z}$), thus the generated cuts
do not separate the solution from the feasible search region for the algorithm.
The reason for that to happen is the combination of the coefficients
obtained by solving the CGLP. As explained in Section \ref{sec:imp3-disjcut},
fractional LP solution of RMP $\bar{z}$, projected on to the variable space
based on the links of the graph and the vehicles. A disjunctive cut which is
violating the current projected fractional solution is found first (by solving
the CGLP), but before projecting the obtained cut to the master problem
space, the coefficients are combined accordingly to get rid of the vehicle
index. Therefore, the new cut is not necessarily violates the
current solution in either space. Therefore, implementation 3 performs very
poorly with respect to LP bound improvement as well as the time. Therefore,
we can easily eliminate implementation 3 of the disjunctive procedure
from the set of cuts considered to extend the algorithms.

{\em Implementation 4} was designed to overcome the implementation 3's problem of
generating cuts which do not violate the LP solution. Instead of combining
the coefficients after the cut is generated, implementation 4 adds new conditions
to the CGLP to eliminate the combination procedure. Thus, the generated cuts
with implementation 4 provide coefficients without vehicle index, i.e.,
the vehicle index is dropped at the beginning because of equality conditions.
Therefore, the generated cut must violate the current fractional solution.
Implementation 4 is successful for small size instances, however, our computational
tests with 25-customer instances showed that addition of new conditions to the
CGLP increases the solution time significantly as seen from  Table \ref{tab:disj-time}.
In addition to this drawback, by investigating the LP relaxation solution after
and before the cut, we found that there exists some nodes in which even though the
generated cut improves the LP relaxation value at first, after the additional column
generation steps designed to consider the new cut, the optimal LP solution moves to a
point which has the objective value that is same or only slightly better than the
previous objective value. As seen from the form of the generated cut (\ref{disj_cut_imp3}),
there are coefficients corresponds to directed links between nodes in the problem. Since
the distance matrix is symmetric, the column generation algorithm generates
new columns which include different orientations of routes that are already
included in the previous columns in the previous LP solution. The new LP solution
satisfies the cut, but has the same objective value as before because of these
new equivalent columns. Therefore, the LP solution value does not change while
the set of non zero pairing variables changes. We provide an example to explain this
problem.
\begin{example}
Branch,price and cut algorithm with cuts from DISJCUTGEN(4) is run for instance
p01-f25-v1t1. In one of the nodes, the LP relaxation value before the addition of
the generated cut is equal to the LP relaxation value after the cut is added and
column generation is done. We compared two LP solutions and only listed the non zero
variables which are different in these two solutions.

\noindent{\bf Before cut is added:}
\begin{optprog*}
& \bar{z}_1 & = & 0.188 & : & \mbox{f$_0$}-4-11-\mbox{f$_0$}-22-6-23-\mbox{f$_0$}-13-24-\mbox{f$_0$}-17-\mbox{f$_0$}\\
& \bar{z}_2 & = & 0.16  & : &  \mbox{f$_0$}-3-16-18-12-\mbox{f$_0$}-22-6-5-\mbox{f$_0$}-13-24-\mbox{f$_0$} \\
& \bar{z}_3 & = & 0.028 & : &  \mbox{f$_0$}- 22-6-23-\mbox{f$_0$}-3-18-12-\mbox{f$_0$}-13-5-\mbox{f$_0$}
\end{optprog*}
{\bf After cut is added:}
\begin{optprog*}
& \bar{z}_2 & = & 0.131 & : &  \mbox{f$_0$}-3-16-18-12-\mbox{f$_0$}-22-6-5-\mbox{f$_0$}-13-24-\mbox{f$_0$} \\
& \bar{z}_4 & = & 0.188 & : &  \mbox{f$_0$}-11-4-\mbox{f$_0$}-23-6-22-\mbox{f$_0$}-24-13-\mbox{f$_0$}-17-\mbox{f$_0$} \\
& \bar{z}_5 & = & 0.029  & : &  \mbox{f$_0$}-12-18-16-3-\mbox{f$_0$}-5-6-22-\mbox{f$_0$}-24-13-\mbox{f$_0$} \\
& \bar{z}_6 & = & 0.028 & : &  \mbox{f$_0$}-23-6-22-\mbox{f$_0$}-12-18-3-\mbox{f$_0$}-13-5-\mbox{f$_0$}
\end{optprog*}
After the cut was added, the column generation generated $\bar{z}_4$ and $\bar{z}_6$ and replaced
$\bar{z}_1$  and $\bar{z}_3$ with them. As the sequences of customers in the corresponding
pairings are investigated, equivalence of $\bar{z}_1$ and $\bar{z}_4$, and  $\bar{z}_3$
and $\bar{z}_6$ are obvious. In addition,  $\bar{z}_5$ which is equal to $\bar{z}_2$ was generated
and took nonzero value in the solution.
\end{example}

Even though there are some nodes in which the disjunctive cut from DISJCUTGEN(4) strictly
improves the LP solution value, the nodes which react the cuts in a similar way to the example
reduce the performance of the overall algorithm especially because of the large
CGLP solution times. Therefore, we can conclude that the branch and price algorithm
without DISJCUTGEN(4) is better of than the  the branch, price and cut algorithm
with DISJCUTGEN(4).
